\documentclass[12pt,a4paper]{article}
\usepackage[margin=25mm, headheight=15pt, headsep=10pt]{geometry}
\usepackage{fontspec}
\usepackage[table]{xcolor} % Note: table option is required for rowcolors
\usepackage{titlesec}
\usepackage{tocloft}
\usepackage{fancyhdr}
\usepackage{booktabs}
\usepackage{tcolorbox}
\tcbuselibrary{skins, breakable}
\usepackage[colorlinks=true, linkcolor=emerald, urlcolor=emerald, bookmarks=true]{hyperref}

\definecolor{emerald}{RGB}{0,102,51}
\definecolor{darkred}{RGB}{139,0,0}
\definecolor{lightgray}{RGB}{245,245,245}

\usepackage[nil]{babel}
\babelprovide[import, main]{bengali}
\babelprovide[import]{arabic}
\babelfont{rm}[Renderer=HarfBuzz,Scale=1.0]{Noto Serif Bengali}
\babelfont{sf}[Renderer=HarfBuzz]{Noto Sans Bengali}
\babelfont{tt}[Renderer=HarfBuzz]{Noto Sans Bengali}
\babelfont[arabic]{rm}[Renderer=HarfBuzz,Scale=1.1]{Noto Naskh Arabic}

% Section styling
\titleformat{\section}{\Large\bfseries\color{emerald}}{\thesection.}{0.5em}{}
\titleformat{\subsection}{\large\bfseries\color{emerald!85!black}}{\thesubsection.}{0.5em}{}
\titleformat{\subsubsection}{\normalsize\bfseries\color{emerald!70!black}}{\thesubsubsection.}{0.5em}{}
\titlespacing*{\section}{0pt}{18pt}{8pt}

% Fancy Header/Footer setup
\pagestyle{fancy}
\fancyhf{} % clear all header and footer fields
\fancyhead[L]{\color{emerald}\small\textbf{\nouppercase{\leftmark}}}
\fancyfoot[L]{\scriptsize\color{black!50}{\lat{AI}}-সংকলিত, আলেম-যাচাইকৃত নয়}
\fancyfoot[C]{\color{emerald!70!black}\thepage}
\renewcommand{\headrulewidth}{0.4pt}
\renewcommand{\headrule}{\hbox to\headwidth{\color{emerald}\leaders\hrule height \headrulewidth\hfill}}
\renewcommand{\footrulewidth}{0pt}

% Custom boxes
% 1. Ayat/Hadith Box
\newtcolorbox{ayatbox}[1][]{
    enhanced, breakable, 
    colback=emerald!5, 
    colframe=emerald, 
    boxrule=0.5pt, 
    arc=3pt, 
    left=10pt, right=10pt, top=10pt, bottom=10pt,
    #1
}

% 2. Quote/Translation Box
\newtcolorbox{quotebox}[1][]{
    enhanced, breakable,
    frame hidden,
    colback=lightgray,
    borderline west={3pt}{0pt}{emerald},
    left=15pt, right=10pt, top=10pt, bottom=10pt,
    sharp corners,
    #1
}

% 3. AI-generated notice box
\newtcolorbox{warnbox}[1][]{
    enhanced, breakable,
    colback=red!4,
    colframe=darkred,
    boxrule=0.8pt,
    arc=3pt,
    left=10pt, right=10pt, top=10pt, bottom=10pt,
    #1
}

% Commands
\newcommand{\arab}[1]{\foreignlanguage{arabic}{#1}}
\newcommand{\kib}[1]{\textcolor{emerald}{\textbf{#1}}}

% Latin (English) fallback: Noto Serif Bengali has NO Latin glyphs
\newfontfamily\latfont[Renderer=HarfBuzz]{Noto Serif}
\newcommand{\lat}[1]{{\latfont #1}}

\setlength{\parskip}{5pt}
\setlength{\parindent}{0pt}

% Fix for missing bullet in Noto Serif Bengali
\renewcommand{\labelitemi}{$\bullet$}



\begin{document}
\begin{center}{\Large\bfseries\color{emerald} আল্লাহর রহমত ও তাওবা — কুরআনের আয়াত}\end{center}
\vspace{1em}
\section*{কুরআনের আয়াত — আল্লাহর রহমত, ক্ষমা ও তাওবা}

\begin{warnbox}
 \textbf{গুরুত্বপূর্ণ নোটিশ (\lat{Important} \lat{Notice}):} এই কন্টেন্টটি \lat{AI} (কৃত্রিম বুদ্ধিমত্তা) দিয়ে প্রস্তুত ও সংকলিত। \lat{AI} কঠোর সূত্র-মান নীতি মেনে শুধু নির্ভরযোগ্য উৎস থেকে তথ্য সংগ্রহ করেছে — কেবল সহীহ/হাসান হাদিস ও সহীহ সনদের আসার রাখা হয়েছে, দুর্বল সনদ বাদ দেওয়া হয়েছে। তবে এটি কোনো আলেম বা মুহাদ্দিস কর্তৃক যাচাইকৃত নয়।
\end{warnbox}

\begin{quote}
\textbf{পড়ার নিয়ম:} আয়াতগুলো ৮টি থিমে সাজানো। প্রতিটিতে — সূত্র কার্ড $\rightarrow$ মূল আরবি (যাচাইকৃত) $\rightarrow$ উচ্চারণ $\rightarrow$ বাংলা অর্থ (\lat{pure}) $\rightarrow$ শব্দের ব্যাখ্যা $\rightarrow$ সংক্ষিপ্ত তাফসীর (ইবনে কাসীর, সা'দী, মাআরিফুল কুরআন প্রভৃতি থেকে; পৃষ্ঠা নম্বর ইচ্ছাকৃতভাবে দেওয়া হয়নি — উৎস গ্রন্থ নাম উল্লেখ করা হলো) $\rightarrow$ সংশ্লিষ্ট দলিল। \\
\textbf{মূলনীতি:} শুধু সহীহ/হাসান হাদিস; দুর্বল সনদ বাদ।
\end{quote}

\medskip\hrule\medskip

\subsection*{থিম ১ — আল্লাহর রহমত কত বড় (\lat{The} \lat{Extent} \lat{of} \lat{Mercy})}

\subsubsection*{আয়াত ১ — সূরা আল-আ'রাফ (৭) — আয়াত ১৫৬}

\subsubsection*{১. সূত্র কার্ড}

\begin{table}[h]\centering\small
\begin{tabular}{ll}
বিষয় & বিবরণ \\
\hline
\textbf{সূরা} & আল-আ'রাফ (৭) — মক্কী \\
\textbf{আয়াত} & ১৫৬ \\
\textbf{মূল বিষয়} & আল্লাহর রহমত সবকিছুকে ঘিরে আছে — কোনো গুনাহ রহমতের বাইরে নয় \\
\textbf{প্রসঙ্গ} & মুসা (আঃ)-এর দোয়ার জবাবে আল্লাহ নিজের গুণ বর্ণনা করছেন \\
\hline
\end{tabular}
\end{table}

\subsubsection*{২. মূল আরবি}

\begin{center}\arab{وَرَحْمَتِي وَسِعَتْ كُلَّ شَيْءٍ ۚ فَسَأَكْتُبُهَا لِلَّذِينَ يَتَّقُونَ وَيُؤْتُونَ الزَّكَاةَ وَالَّذِينَ هُمْ بِآيَاتِنَا يُؤْمِنُونَ}\end{center}
\textbf{উচ্চারণ:} ওয়া রাহমাতী ওয়াসি'আত কুল্লা শাই। ফাসাআকতুবুহা লিল্লাযীনা ইয়াত্তাকূনা ওয়া ইউ'তূনায যাকাতা ওয়াল্লাযীনা হুম বিআয়াতিনা ইউ'মিনূন।

\subsubsection*{৩. বাংলা অর্থ (\lat{pure})}

\textbf{\lat{“}আর আমার রহমত প্রতিটি জিনিসকে পরিবেষ্টন করে আছে। সুতরাং আমি তা লিখে দেব তাদের জন্য যারা তাকওয়া অবলম্বন করে, যাকাত দেয় এবং যারা আমার আয়াতসমূহে ঈমান আনে।\lat{”}}

\subsubsection*{৪. শব্দের ব্যাখ্যা}

\begin{table}[h]\centering\small
\begin{tabular}{ll}
শব্দ & অর্থ \\
\hline
\textbf{ওয়াসি'আত (\arab{وسعت})} & পরিবেষ্টন/ঘিরে রাখা — ব্যাপকতা \\
\textbf{কুল্লা শাই (\arab{كل} \arab{شيء})} & প্রতিটি জিনিস — সীমাহীন ব্যাপ্তি \\
\textbf{ফাসাআকতুবুহা} & সুতরাং আমি তা লিখে দেব \\
\hline
\end{tabular}
\end{table}

\subsubsection*{৫. সংক্ষিপ্ত তাফসীর}

\begin{itemize}
\item \textbf{ইবনে কাসীর:} রহমত বলতে এখানে আল্লাহর দয়া ও ক্ষমার ব্যাপকতা — যা আসমান, জমিন, জিন-ইনসান সবকিছুকে ঘিরে। মুসা (আঃ) উম্মতের জন্য রহমত চাইলে আল্লাহ বলেন — আমার রহমত সবকিছুর চেয়ে বড়; কিন্তু দুনিয়ায় তার বিশেষ লেখা (অধিকাংশ রহমত) তাদের জন্য, যারা তাকওয়া ও ঈমান রাখে। অর্থ — রহমত সবার জন্য উন্মুক্ত, তবে পূর্ণতা তাকওয়ার সাথে।
\item \textbf{সা'দী:} \lat{“}রহমত সবকিছুকে ঘিরে আছে\lat{”} — এর মধ্যে তাওবাকারী গুনাহগারও আছে; তবে যারা আল্লাহর পথে এগিয়ে থাকে, তাদের জন্য রহমতের বিশেষ ব্যবস্থা। এই আয়াতই বলে — পাপীর আশা কখনো মরতে পারে না, রহমত তাকে ঘিরে আছে; তাকে শুধু ফিরে আসতে হবে।
\item \textbf{মাআরিফুল কুরআন:} এই আয়াতের শিক্ষা — আল্লাহর রহমত তাঁর গজবের চেয়ে ব্যাপক; তাই \lat{“}আমার পাপ অনেক, আল্লাহ ক্ষমা করবেন না\lat{”} — এই ধারণা কুরআনবিরোধী। রহমতের দরজা কেবল হতাশার জন্যই বন্ধ, তাওবার জন্য নয়।
\end{itemize}
\subsubsection*{৬. সংশ্লিষ্ট দলিল}

\begin{itemize}
\item \textbf{আয়াত (যুমার ৩৯:৫৩):} \lat{“}তোমরা আল্লাহর রহমত থেকে নিরাশ হয়ো না\lat{”} — একই বার্তার ঘোষণা।
\item \textbf{হাদিস (মুসলিম ২৭৫১):} \lat{“}নিশ্চয়ই আমার রহমত আমার গজবের উপর প্রাধান্য পায়।\lat{”}
\end{itemize}
\medskip\hrule\medskip

\subsubsection*{আয়াত ২ — সূরা আল-আন'আম (৬) — আয়াত ১২ ও ৫৪}

\subsubsection*{১. সূত্র কার্ড}

\begin{table}[h]\centering\small
\begin{tabular}{ll}
বিষয় & বিবরণ \\
\hline
\textbf{সূরা} & আল-আন'আম (৬) — মক্কী \\
\textbf{আয়াত} & ১২ ও ৫৪ \\
\textbf{মূল বিষয়} & আল্লাহ নিজের উপর রহমত লিখে নিয়েছেন — নিজেকে রহমতের জন্য দায়ী করেছেন \\
\textbf{প্রসঙ্গ} & মুশরিকদের প্রশ্নের উত্তর; মুমিনদের কাছে ফিরে আসা পাপীদের স্বাগত \\
\hline
\end{tabular}
\end{table}

\subsubsection*{২. মূল আরবি}

\begin{center}\arab{قُلْ لِمَنْ مَا فِي السَّمَاوَاتِ وَالْأَرْضِ ۖ قُلْ لِلَّهِ ۚ كَتَبَ عَلَىٰ نَفْسِهِ الرَّحْمَةَ ۚ لَيَجْمَعَنَّكُمْ إِلَىٰ يَوْمِ الْقِيَامَةِ لَا رَيْبَ فِيهِ}\end{center}
\textbf{উচ্চারণ:} কুল লিমান মা-ফিস সামাওয়াতি ওয়াল আরদ। কুল লিল্লাহ। কাতাবা 'আলা- নাফসিহির রাহমাহ। লাইয়াজমা'আন্নাকুম ইলা- ইয়াওমিল কিয়ামাতি লা- রাইবা ফীহ।

\begin{center}\arab{وَإِذَا جَاءَكَ الَّذِينَ يُؤْمِنُونَ بِآيَاتِنَا فَقُلْ سَلَامٌ عَلَيْكُمْ ۖ كَتَبَ رَبُّكُمْ عَلَىٰ نَفْسِهِ الرَّحْمَةَ ۖ أَنَّهُ مَنْ عَمِلَ مِنْكُمْ سُوءًا بِجَهَالَةٍ ثُمَّ تَابَ مِنْ بَعْدِهِ وَأَصْلَحَ فَأَنَّهُ غَفُورٌ رَحِيمٌ}\end{center}
\textbf{উচ্চারণ:} ওয়া ইযা- জাআকাল্লাযীনা ইউ'মিনূনা বিআয়াতিনা ফাকুল সালামুন 'আলাইকুম। কাতাবা রব্বুকুম 'আলা- নাফসিহির রাহমাহ, আন্নাহু মান 'আমিলা মিনকুম সূআন বিজাহালাতিন ছুম্মা তা-বা মিম বা'দিহী ওয়া আসলাহা ফাআন্নাহু গাফূরুর রাহীম।

\subsubsection*{৩. বাংলা অর্থ (\lat{pure})}

\textbf{\lat{“}বলো: আসমান ও জমিনে যা আছে, কার? বলো: আল্লাহর। তিনি নিজের উপর রহমত লিখে নিয়েছেন। তিনি অবশ্যই তোমাদেরকে কিয়ামতের দিন একত্র করবেন — এতে কোনো সন্দেহ নেই।\lat{”}} (৬:১২)

\textbf{\lat{“}আর যখন তোমার কাছে তারা আসে, যারা আমার আয়াতে ঈমান আনে, তখন তুমি বলো: তোমাদের উপর সালাম। তোমাদের রব নিজের উপর রহমত লিখে নিয়েছেন — তোমাদের মধ্যে কেউ অজ্ঞতাবশত মন্দ কাজ করল, এরপর তাওবা করল এবং সংশোধন হয়ে গেল, তবে তিনি (জেনে রাখো) ক্ষমাশীল, পরম দয়ালু।\lat{”}} (৬:৫৪)

\subsubsection*{৪. শব্দের ব্যাখ্যা}

\begin{table}[h]\centering\small
\begin{tabular}{ll}
শব্দ & অর্থ \\
\hline
\textbf{কাতাবা 'আলা নাফসিহির রাহমাহ (\arab{كتب} \arab{على} \arab{نفسه} \arab{الرحمة})} & নিজের উপর রহমত লেখা — রহমতকে নিজের দায়িত্ব বানানো \\
\textbf{বিজাহালাতিন (\arab{بجهالة})} & অজ্ঞতাবশত — এখানে অজ্ঞতা মানে গুনাহের মাত্রা না বোঝা, শিরক নয় \\
\textbf{থুম্মা তাবা} & এরপর তাওবা করল \\
\textbf{ওয়া আসলাহা (\arab{وأصلح})} & এবং সংশোধন/সংস্কার করল — তাওবার পরের অবস্থা \\
\hline
\end{tabular}
\end{table}

\subsubsection*{৫. সংক্ষিপ্ত তাফসীর}

\begin{itemize}
\item \textbf{ইবনে কাসীর:} \lat{“}নিজের উপর রহমত লিখে নিয়েছেন\lat{”} — এটা আল্লাহর পক্ষ থেকে রহমতের প্রতিশ্রুতি; ইবনে আব্বাস বলেন, এটি এমন একটি বাক্য যা আল্লাহ নিজের উপর ওয়াজিব করেছেন। ৬:৫৪-এর ব্যাখ্যায় ইবনে কাসীর বলেন — এই আয়াত নাজিল হয়েছিল এমন একদল লোকের ব্যাপারে যারা ইসলামের আগে বড় বড় পাপ করেছিল এবং ভাবত আল্লাহ তাদের ক্ষমা করবেন না; আল্লাহ বললেন — \lat{“}তোমাদের উপর সালাম\lat{”}, তাওবা করলে ক্ষমা আছে।
\item \textbf{সা'দী:} ৬:৫৪ হচ্ছে পাপীদের জন্য কুরআনের সবচেয়ে বড় সুসংবাদগুলোর একটি — যে ব্যক্তি যেকোনো গুনাহ করে ফেলে (শিরক ব্যতীত) এবং আন্তরিক তাওবা করে, আল্লাহ তার ক্ষমা প্রতিশ্রুত করেছেন।
\item \textbf{মাআরিফুল কুরআন:} ৬:১২-এর আলোকে — রহমতের এত ব্যাপকতা সত্ত্বেও মানুষ যদি কুফরি অবস্থায় মরে, তবে রহমতের প্রতিশ্রুতি তাকে ঢাকবে না; রহমত পেতে হলে ফিরে আসতে হবে।
\end{itemize}
\subsubsection*{৬. সংশ্লিষ্ট দলিল}

\begin{itemize}
\item \textbf{হাদিস (মুসলিম ২৭৫১):} \lat{“}আমার রহমত আমার গজবের উপর প্রাধান্য পায়\lat{”} — একই অর্থ।
\item \textbf{আয়াত (আ'রাফ ৭:১৫৬):} \lat{“}আমার রহমত সবকিছুকে ঘিরে আছে।\lat{”}
\end{itemize}
\medskip\hrule\medskip

\subsubsection*{আয়াত ৩ — সূরা আল-মু'মিন (গাফির) (৪০) — আয়াত ৩}

\subsubsection*{১. সূত্র কার্ড}

\begin{table}[h]\centering\small
\begin{tabular}{ll}
বিষয় & বিবরণ \\
\hline
\textbf{সূরা} & আল-মু'মিন (গাফির) (৪০) — মক্কী \\
\textbf{আয়াত} & ৩ \\
\textbf{মূল বিষয়} & আল্লাহর পরিচয় — গুনাহ ক্ষমাকারী ও তাওবা গ্রহণকারী \\
\textbf{প্রসঙ্গ} & কুরআনের অবতরণের পর আল্লাহর পরিচয়ের ঘোষণা \\
\hline
\end{tabular}
\end{table}

\subsubsection*{২. মূল আরবি}

\begin{center}\arab{غَافِرِ الذَّنْبِ وَقَابِلِ التَّوْبِ شَدِيدِ الْعِقَابِ ذِي الطَّوْلِ ۖ لَا إِلَٰهَ إِلَّا هُوَ ۖ إِلَيْهِ الْمَصِيرُ}\end{center}
\textbf{উচ্চারণ:} গা-ফিরিয যাম্বি ওয়া কা-বিলিত তাওবি শাদীদিল 'ইকা-বি যিত্ব তাওল। লা- ইলা-হা ইল্লা- হুওয়া, ইলাইহিল মাসীর।

\subsubsection*{৩. বাংলা অর্থ (\lat{pure})}

\textbf{\lat{“}তিনি গুনাহ ক্ষমাকারী, তাওবা গ্রহণকারী, কঠোর শাস্তিদাতা, অনুগ্রহের মালিক। তিনি ছাড়া কোনো সত্য ইলাহ নেই — তাঁরই দিকে প্রত্যাবর্তন।\lat{”}}

\subsubsection*{৪. শব্দের ব্যাখ্যা}

\begin{table}[h]\centering\small
\begin{tabular}{ll}
শব্দ & অর্থ \\
\hline
\textbf{গাফিরিয যাম্ব (\arab{غافر} \arab{الذنب})} & গুনাহ ক্ষমাকারী — গুনাহকে ঢেকে দেওয়া \\
\textbf{কাবিলুত তাওব (\arab{قابل} \arab{التوب})} & তাওবা গ্রহণকারী — তাওবার দরজা খোলা \\
\textbf{শাদীদিল ইকাব} & কঠোর শাস্তিদাতা — ভয়ের ভারসাম্য \\
\textbf{যিত্ব তাওল (\arab{ذي} \arab{الطول})} & অনুগ্রহ/দানশীলতার মালিক \\
\hline
\end{tabular}
\end{table}

\subsubsection*{৫. সংক্ষিপ্ত তাফসীর}

\begin{itemize}
\item \textbf{ইবনে কাসীর:} এখানে আল্লাহ তিনটি নাম/গুণ একসাথে উল্লেখ করেছেন — ক্ষমা, তাওবা গ্রহণ ও কঠিন শাস্তি। অর্থ — রহমত ও শাস্তি দুই-ই তাঁর; বান্দার কর্তব্য — ক্ষমার প্রতিশ্রুতিতে আশা রাখা, আবার শাস্তির ভয়ে সাবধান থাকা।
\item \textbf{সা'দী:} \lat{“}তাওবা গ্রহণকারী\lat{”} বলেই বোঝা যায় — যতবার বান্দা তাওবা করবে, আল্লাহ ততবার গ্রহণ করবেন; তাওবা শেষ হয়ে যায় বলে কিছু নেই। সাথে \lat{“}কঠোর শাস্তিদাতা\lat{”} — যেন তাওবাকে হালকা না নেয়।
\item \textbf{মাআরিফুল কুরআন:} এই আয়াতে আশা ও ভয় দুই-ই; পাপী ব্যক্তির জন্য আশার দিকটি হলো — ক্ষমা ও তাওবা গ্রহণের গুণ তাকে ডাকছে, শাস্তির ভয় তাকে সাবধান করছে।
\end{itemize}
\subsubsection*{৬. সংশ্লিষ্ট দলিল}

\begin{itemize}
\item \textbf{আয়াত (যুমার ৩৯:৫৩):} ক্ষমার সর্বোচ্চ ঘোষণা।
\item \textbf{হাদিস (মুসলিম ২৭৫৯):} \lat{“}আল্লাহ রাতে তাঁর হাত বাড়ান দিনের পাপীর তাওবার জন্য…\lat{”}
\end{itemize}
\medskip\hrule\medskip

\subsection*{থিম ২ — রহমত থেকে নিরাশ হওয়া হারাম (\lat{No} \lat{Despair})}

\subsubsection*{আয়াত ৪ — সূরা আয-যুমার (৩৯) — আয়াত ৫৩ (সবচেয়ে গুরুত্বপূর্ণ আয়াত)}

\subsubsection*{১. সূত্র কার্ড}

\begin{table}[h]\centering\small
\begin{tabular}{ll}
বিষয় & বিবরণ \\
\hline
\textbf{সূরা} & আয-যুমার (৩৯) — মক্কী \\
\textbf{আয়াত} & ৫৩ \\
\textbf{মূল বিষয়} & পাপে সীমালঙ্ঘনকারীদের জন্য ক্ষমার সরাসরি ঘোষণা — হতাশা নিষিদ্ধ \\
\textbf{প্রসঙ্গ} & এই আয়াতের ব্যাপারে সাহাবিরা বলেছেন — কুরআনের সবচেয়ে আশাপূর্ণ আয়াত \\
\hline
\end{tabular}
\end{table}

\subsubsection*{২. মূল আরবি}

\begin{center}\arab{قُلْ يَا عِبَادِيَ الَّذِينَ أَسْرَفُوا عَلَىٰ أَنْفُسِهِمْ لَا تَقْنَطُوا مِنْ رَحْمَةِ اللَّهِ ۚ إِنَّ اللَّهَ يَغْفِرُ الذُّنُوبَ جَمِيعًا ۚ إِنَّهُ هُوَ الْغَفُورُ الرَّحِيمُ}\end{center}
\textbf{উচ্চারণ:} কুল ইয়া- 'ইবা-দিয়াল্লাযীনা আসরাফূ 'আলা- আনফুসিহিম লা- তাকনাতূ মির রাহমাতিল্লাহ। ইন্নাল্লা-হা ইয়াগফিরুজ যুনূবা জামী'আ। ইন্নাহু হুওয়াল গাফূরুর রাহীম।

\subsubsection*{৩. বাংলা অর্থ (\lat{pure})}

\textbf{\lat{“}বলো: হে আমার বান্দারা! যারা নিজেদের উপর সীমালঙ্ঘন (জুলুম) করেছ, তোমরা আল্লাহর রহমত থেকে নিরাশ হয়ো না। নিশ্চয়ই আল্লাহ সমস্ত গুনাহ ক্ষমা করে দেন। নিশ্চয়ই তিনি ক্ষমাশীল, পরম দয়ালু।\lat{”}}

\subsubsection*{৪. শব্দের ব্যাখ্যা}

\begin{table}[h]\centering\small
\begin{tabular}{ll}
শব্দ & অর্থ \\
\hline
\textbf{আসরাফূ (\arab{أسرفوا})} & সীমালঙ্ঘন করা — গুনাহতে বাড়াবাড়ি করা \\
\textbf{লা তাকনাতূ (\arab{لا} \arab{تقنطوا})} & নিরাশ হয়ো না — হতাশা সম্পূর্ণ নিষিদ্ধ \\
\textbf{ইয়াগফিরুয যুনূবা জামী'আ} & সমস্ত গুনাহ ক্ষমা করে দেন — ব্যাপক ঘোষণা \\
\hline
\end{tabular}
\end{table}

\subsubsection*{৫. সংক্ষিপ্ত তাফসীর}

\begin{itemize}
\item \textbf{ইবনে কাসীর:} আলী, ইবনে মাসউদ, ইবনে আব্বাস, আবু সাঈদ খুদরি (রাঃ)-সহ বহু সাহাবি বলেছেন — কুরআনে সবচেয়ে আশাপূর্ণ (\lat{most} \lat{hopeful}) আয়াত এটি। এখানে \lat{“}আসরাফূ\lat{”} শব্দটি মুশরিকসহ সকল গুনাহগারকে ডাকছে — ইসলাম গ্রহণের আগে যারা শিরকও করেছে, তারা তাওবার মাধ্যমে ক্ষমা পেতে পারে; শিরকের উপর মৃত্যু হলে তা নয় (নিসা ৪:৪৮)। তাফসীরকারকগণ বলেন — এই আয়াত সবচেয়ে বড় \lat{“}প্রলোভন\lat{”}-এর জবাব: গুনাহ যত বড়ই হোক, রহমত থেকে নিরাশ হওয়া আল্লাহর অপবাদ।
\item \textbf{সা'দী:} এটি রহমতের সর্বোচ্চ ঘোষণা — আল্লাহ নিজেই পাপীদের ডেকে বলছেন \lat{“}হে আমার বান্দারা\lat{”} — পরিচয়ের স্নেহে; নির্দেশ করছেন হতাশ না হতে। কেননা হতাশা দুটো জিনিস করে — (১) বান্দা তাওবা ছেড়ে দেয়, (২) আল্লাহর ক্ষমা-রহমত সম্পর্কে খারাপ ধারণা পোষণ করে।
\item \textbf{মাআরিফুল কুরআন:} এ আয়াত নাজিলের প্রেক্ষাপটে উল্লেখ আছে — একদল সাহাবি ইসলামের আগের পাপের জন্য চিন্তিত ছিলেন; এই আয়াত তাদের আশ্বস্ত করল। শিক্ষা — অতীত যতই ভারী হোক, ইসলামের দরজা খোলা; \lat{“}আল্লাহ আমাকে ক্ষমা করবেন না\lat{”} — এই চিন্তা শয়তানেরই কৌশল।
\end{itemize}
\subsubsection*{৬. সংশ্লিষ্ট দলিল}

\begin{itemize}
\item \textbf{আয়াত (ইউসুফ ১২:৮৭):} \lat{“}আল্লাহর রহমত থেকে শুধু কাফির সম্প্রদায়ই নিরাশ হয়।\lat{”}
\item \textbf{হাদিস (তিরমিযী ৩৫৪০):} \lat{“}তোমার গুনাহ আকাশের মেঘ পর্যন্ত পৌঁছালেও আমি ক্ষমা করব।\lat{”}
\item \textbf{আয়াত (নিসা ৪:৪৮):} শিরকের উপর মৃত্যুর ক্ষেত্রে ক্ষমা নেই — সীমারেখা।
\end{itemize}
\medskip\hrule\medskip

\subsubsection*{আয়াত ৫ — সূরা ইউসুফ (১২) — আয়াত ৮৭}

\subsubsection*{১. সূত্র কার্ড}

\begin{table}[h]\centering\small
\begin{tabular}{ll}
বিষয় & বিবরণ \\
\hline
\textbf{সূরা} & ইউসুফ (১২) — মক্কী \\
\textbf{আয়াত} & ৮৭ \\
\textbf{মূল বিষয়} & রহমত থেকে নিরাশ হওয়া কাফিরদের বৈশিষ্ট্য — মুমিনের নয় \\
\textbf{প্রসঙ্গ} & ইয়াকুব (আঃ) পুত্রদের ইউসুফের খোঁজে পাঠানোর সময় \\
\hline
\end{tabular}
\end{table}

\subsubsection*{২. মূল আরবি}

\begin{center}\arab{يَا بَنِيَّ اذْهَبُوا فَتَحَسَّسُوا مِنْ يُوسُفَ وَأَخِيهِ وَلَا تَيْأَسُوا مِنْ رَوْحِ اللَّهِ ۖ إِنَّهُ لَا يَيْأَسُ مِنْ رَوْحِ اللَّهِ إِلَّا الْقَوْمُ الْكَافِرُونَ}\end{center}
\textbf{উচ্চারণ:} ইয়া- বানিয়্যায হাবূ ফাতাহাস্সাসূ মিই ইউসুফা ওয়া আখীহি ওয়ালা- তাইয়াসূ মির রাওহিল্লাহ। ইন্নাহু লা- ইয়াইয়াসু মির রাওহিল্লাহি ইল্লাল কাওমুল কা-ফিরূন।

\subsubsection*{৩. বাংলা অর্থ (\lat{pure})}

\textbf{\lat{“}হে আমার ছেলেরা! তোমরা যাও এবং ইউসুফ ও তার ভাইয়ের খোঁজ করো। আর আল্লাহর রহমত (সান্ত্বনা) থেকে নিরাশ হয়ো না। নিশ্চয়ই আল্লাহর রহমত থেকে কেবল কাফির সম্প্রদায়ই নিরাশ হয়।\lat{”}}

\subsubsection*{৪. শব্দের ব্যাখ্যা}

\begin{table}[h]\centering\small
\begin{tabular}{ll}
শব্দ & অর্থ \\
\hline
\textbf{তাহাস্সাসূ (\arab{تحسسوا})} & খোঁজ করা — অনুসন্ধান চালানো \\
\textbf{রাওহুল্লাহ (\arab{روح} \arab{الله})} & আল্লাহর রহমত/সান্ত্বনা — উদ্ধারের সুবাতাস \\
\textbf{লা ইয়াইয়াসু} & নিরাশ হয় না \\
\hline
\end{tabular}
\end{table}

\subsubsection*{৫. সংক্ষিপ্ত তাফসীর}

\begin{itemize}
\item \textbf{ইবনে কাসীর:} দীর্ঘ বিচ্ছেদের পরও ইয়াকুব (আঃ) নিরাশ হননি — উল্টো সন্তানদের বললেন খোঁজ চালাতে, কেননা আল্লাহর রহমত থেকে নিরাশ হওয়া কাফিরদের কাজ। এই আয়াত থেকে শিক্ষা — বিপদ ও গুনাহের অন্ধকারে আশা রাখাই ঈমানের পরিচয়।
\item \textbf{সা'দী:} \lat{“}রওহুল্লাহ\lat{”} অর্থ — আল্লাহর রহমত, উদ্ধার ও ত্রাণ। যে ব্যক্তি জানে আল্লাহ তাঁর উপর রহমত করতে পারেন, সে কখনো নিরাশ হবে না; নিরাশা আল্লাহর ক্ষমতা সম্পর্কে খারাপ ধারণা — যা কুফরি পর্যায়ের।
\item \textbf{মাআরিফুল কুরআন:} পাপী ব্যক্তির জন্যও একই শিক্ষা — নিজের অতীত যত খারাপই হোক, \lat{“}ফিরে আসার\lat{”} আশা কখনো ছাড়া যাবে না; কারণ সেই আশাই তাকে আল্লাহর দিকে টানে।
\end{itemize}
\subsubsection*{৬. সংশ্লিষ্ট দলিল}

\begin{itemize}
\item \textbf{আয়াত (হিজর ১৫:৫৬):} \lat{“}নিজের রবের রহমত থেকে কে নিরাশ হয় — পথভ্রষ্টরা ছাড়া?\lat{”}
\item \textbf{আয়াত (যুমার ৩৯:৫৩):} ক্ষমার সর্বোচ্চ ঘোষণা।
\end{itemize}
\medskip\hrule\medskip

\subsubsection*{আয়াত ৬ — সূরা আল-হিজর (১৫) — আয়াত ৫৬}

\subsubsection*{১. সূত্র কার্ড}

\begin{table}[h]\centering\small
\begin{tabular}{ll}
বিষয় & বিবরণ \\
\hline
\textbf{সূরা} & আল-হিজর (১৫) — মক্কী \\
\textbf{আয়াত} & ৫৬ \\
\textbf{মূল বিষয়} & রহমত থেকে নিরাশ হওয়া পথভ্রষ্টদের বৈশিষ্ট্য \\
\textbf{প্রসঙ্গ} & ইবরাহীম (আঃ)-এর সন্তান-সংবাদ পাওয়ার প্রতিক্রিয়া \\
\hline
\end{tabular}
\end{table}

\subsubsection*{২. মূল আরবি}

\begin{center}\arab{قَالَ وَمَنْ يَقْنَطُ مِنْ رَحْمَةِ رَبِّهِ إِلَّا الضَّالُّونَ}\end{center}
\textbf{উচ্চারণ:} কা-লা ওয়া মাই ইয়াকনাতু মির রাহমাতি রব্বিহী ইল্লাদ দা-ল্লূন।

\subsubsection*{৩. বাংলা অর্থ (\lat{pure})}

\textbf{\lat{“}সে (ইবরাহীম) বলল: নিজের রবের রহমত থেকে কে নিরাশ হয় — পথভ্রষ্টরা ছাড়া!\lat{”}}

\subsubsection*{৪. শব্দের ব্যাখ্যা}

\begin{table}[h]\centering\small
\begin{tabular}{ll}
শব্দ & অর্থ \\
\hline
\textbf{ইয়াকনাতু (\arab{يقنط})} & নিরাশ হয় \\
\textbf{আদ-দাল্লূন (\arab{الضالون})} & পথভ্রষ্টরা — যারা আল্লাহ সম্পর্কে ভুল ধারণা পোষণ করে \\
\hline
\end{tabular}
\end{table}

\subsubsection*{৫. সংক্ষিপ্ত তাফসীর}

\begin{itemize}
\item \textbf{ইবনে কাসীর:} বৃদ্ধ বয়সে সন্তানের সুসংবাদ পেয়ে ইবরাহীম (আঃ) বিস্মিত হলেন, কিন্তু প্রশ্ন করলেন না \lat{“}আমি কীভাবে সন্তান পাব\lat{”} — বরং বললেন রহমত থেকে কেউ নিরাশ হয় না। এখানে শিক্ষা — নবীগণও রহমতের আশা কখনো ছাড়েননি।
\item \textbf{সা'দী:} নিরাশা পথভ্রষ্টতার আলামত — কারণ যে আল্লাহকে জানে, সে জানে তাঁর হাতে সবকিছু; তাই নিরাশ হওয়া মানে আল্লাহর ক্ষমতা সম্পর্কে সন্দেহ।
\end{itemize}
\subsubsection*{৬. সংশ্লিষ্ট দলিল}

\begin{itemize}
\item \textbf{আয়াত (যুমার ৩৯:৫৩), (ইউসুফ ১২:৮৭):} নিরাশা নিষেধের তিনটি স্তর।
\end{itemize}
\medskip\hrule\medskip

\subsection*{থিম ৩ — তাওবার দরজা খোলা (\lat{The} \lat{Door} \lat{of} \lat{Tawbah})}

\subsubsection*{আয়াত ৭ — সূরা আশ-শূরা (৪২) — আয়াত ২৫}

\subsubsection*{১. সূত্র কার্ড}

\begin{table}[h]\centering\small
\begin{tabular}{ll}
বিষয় & বিবরণ \\
\hline
\textbf{সূরা} & আশ-শূরা (৪২) — মক্কী \\
\textbf{আয়াত} & ২৫ \\
\textbf{মূল বিষয়} & আল্লাহ নিজেই বান্দাদের তাওবা গ্রহণ করেন \\
\textbf{প্রসঙ্গ} & বান্দাদের প্রতি আল্লাহর দয়ার পরিচয় \\
\hline
\end{tabular}
\end{table}

\subsubsection*{২. মূল আরবি}

\begin{center}\arab{وَهُوَ الَّذِي يَقْبَلُ التَّوْبَةَ عَنْ عِبَادِهِ وَيَعْفُو عَنِ السَّيِّئَاتِ وَيَعْلَمُ مَا تَفْعَلُونَ}\end{center}
\textbf{উচ্চারণ:} ওয়া হুওয়াল্লাযী ইয়াকবালুত তাওবাতা 'আন 'ইবা-দিহী ওয়া ইয়া'ফূ 'আনিস সাইয়ি-আতি ওয়া ইয়া'লামু মা- তাফ'আলূন।

\subsubsection*{৩. বাংলা অর্থ (\lat{pure})}

\textbf{\lat{“}আর তিনিই তাঁর বান্দাদের কাছ থেকে তাওবা গ্রহণ করেন এবং পাপসমূহ ক্ষমা করেন, আর তোমরা যা করো সে সবই তিনি জানেন।\lat{”}}

\subsubsection*{৪. শব্দের ব্যাখ্যা}

\begin{table}[h]\centering\small
\begin{tabular}{ll}
শব্দ & অর্থ \\
\hline
\textbf{ইয়াকবালুত তাওবাহ (\arab{يقبل} \arab{التوبة})} & তাওবা গ্রহণ করেন \\
\textbf{ইয়া'ফূ (\arab{يعفو})} & ক্ষমা/মার্জনা করেন \\
\hline
\end{tabular}
\end{table}

\subsubsection*{৫. সংক্ষিপ্ত তাফসীর}

\begin{itemize}
\item \textbf{ইবনে কাসীর:} এই আয়াত তাওবার গ্রহণকে আল্লাহর নিজের কাজ হিসেবে বর্ণনা করেছে — বান্দার তাওবা কোনো \lat{“}মানুষের দয়ার\lat{”} উপর নির্ভর করে না, বরং আল্লাহ নিজেই গ্রহণ করেন। এটাই সবচেয়ে বড় উৎসাহ — তাওবা করলেই গ্রহণ হবে, কারণ গ্রহণকারী নিজেই ঘোষণা দিয়েছেন।
\item \textbf{সা'দী:} আয়াতের শেষাংশ \lat{“}তোমরা যা করো তিনি জানেন\lat{”} — অর্থ: তাওবা আসল না আসল, তিনি জানেন; তাই তাওবা হতে হবে অন্তর থেকে, শুধু মুখে নয়।
\end{itemize}
\subsubsection*{৬. সংশ্লিষ্ট দলিল}

\begin{itemize}
\item \textbf{আয়াত (তাওবা ৯:১০৪):} \lat{“}তারা কি জানে না — আল্লাহই তাওবা গ্রহণ করেন?\lat{”}
\item \textbf{হাদিস (মুসলিম ২৭৫৯):} রাত-দিন হাত বাড়ানো।
\end{itemize}
\medskip\hrule\medskip

\subsubsection*{আয়াত ৮ — সূরা আত-তাওবা (৯) — আয়াত ১০৪}

\subsubsection*{১. সূত্র কার্ড}

\begin{table}[h]\centering\small
\begin{tabular}{ll}
বিষয় & বিবরণ \\
\hline
\textbf{সূরা} & আত-তাওবা (৯) — মাদানী \\
\textbf{আয়াত} & ১০৪ \\
\textbf{মূল বিষয়} & তাওবা ও সাদাকার গ্রহণকারী আল্লাহ নিজেই \\
\textbf{প্রসঙ্গ} & তাওবাকারীদের উৎসাহ ও সংশয় দূরীকরণ \\
\hline
\end{tabular}
\end{table}

\subsubsection*{২. মূল আরবি}

\begin{center}\arab{أَلَمْ يَعْلَمُوا أَنَّ اللَّهَ هُوَ يَقْبَلُ التَّوْبَةَ عَنْ عِبَادِهِ وَيَأْخُذُ الصَّدَقَاتِ وَأَنَّ اللَّهَ هُوَ التَّوَّابُ الرَّحِيمُ}\end{center}
\textbf{উচ্চারণ:} আলাম ইয়া'লামূ আন্নাল্লা-হা হুওয়া ইয়াকবালুত তাওবাতা 'আন 'ইবা-দিহী ওয়া ইয়া'খুযুস সাদাকা-তি ওয়া আন্নাল্লা-হা হুওয়াত তাওওয়াবুর রাহীম।

\subsubsection*{৩. বাংলা অর্থ (\lat{pure})}

\textbf{\lat{“}তারা কি জানে না — আল্লাহই তাঁর বান্দাদের কাছ থেকে তাওবা গ্রহণ করেন এবং সাদাকা কবুল করেন? আর নিশ্চয়ই আল্লাহ তাওবা গ্রহণকারী, পরম দয়ালু।\lat{”}}

\subsubsection*{৪. শব্দের ব্যাখ্যা}

\begin{table}[h]\centering\small
\begin{tabular}{ll}
শব্দ & অর্থ \\
\hline
\textbf{আত-তাওওয়াব (\arab{التواب})} & তাওবা গ্রহণকারী — বারবার গ্রহণকারী (বহুবার তাওবা গ্রহণের গুণ) \\
\hline
\end{tabular}
\end{table}

\subsubsection*{৫. সংক্ষিপ্ত তাফসীর}

\begin{itemize}
\item \textbf{ইবনে কাসীর:} \lat{“}তাওওয়াব\lat{”} নামটি বহুবার গ্রহণের ইঙ্গিত দেয় — বান্দা যতবার ফিরে আসবে, আল্লাহ ততবার গ্রহণ করবেন। ইবনে কাসীর এই আয়াতের ব্যাখ্যায় বলেছেন — আল্লাহর কাছে তাওবার কোনো \lat{“}সীমা\lat{”} নেই; মানুষের দুর্বলতাই তাওবার কারণ, আল্লাহর রহমতই তার জবাব।
\item \textbf{সা'দী:} এখানে তাওবা ও সাদাকা একসাথে — কারণ দুই-ই \lat{“}ফিরে আসা\lat{”}র কাজ: তাওবা পাপ থেকে, সাদাকা সম্পদ থেকে।
\end{itemize}
\subsubsection*{৬. সংশ্লিষ্ট দলিল}

\begin{itemize}
\item \textbf{আয়াত (শূরা ৪২:২৫):} তাওবা গ্রহণের একই ঘোষণা।
\item \textbf{হাদিস (ইবনে মাজাহ ৪২৫১):} \lat{“}সেরা গুনাহকারী তাওবাকারী।\lat{”}
\end{itemize}
\medskip\hrule\medskip

\subsubsection*{আয়াত ৯ — সূরা আল-বাকারা (২) — আয়াত ১৬০}

\subsubsection*{১. সূত্র কার্ড}

\begin{table}[h]\centering\small
\begin{tabular}{ll}
বিষয় & বিবরণ \\
\hline
\textbf{সূরা} & আল-বাকারা (২) — মাদানী \\
\textbf{আয়াত} & ১৬০ \\
\textbf{মূল বিষয়} & সত্য গোপনকারীদের জন্যও তাওবার দরজা \\
\textbf{প্রসঙ্গ} & পূর্ববর্তী কিতাবের আহলে কিতাব যারা সত্য গোপন করেছিল \\
\hline
\end{tabular}
\end{table}

\subsubsection*{২. মূল আরবি}

\begin{center}\arab{إِلَّا الَّذِينَ تَابُوا وَأَصْلَحُوا وَبَيَّنُوا فَأُولَٰئِكَ أَتُوبُ عَلَيْهِمْ ۚ وَأَنَا التَّوَّابُ الرَّحِيمُ}\end{center}
\textbf{উচ্চারণ:} ইল্লাল্লাযীনা তা-বূ ওয়া আসলাহূ ওয়া বাইয়্যানূ ফাউলা-ইকা আতূবু 'আলাইহিম। ওয়া আনাত তাওওয়াবুর রাহীম।

\subsubsection*{৩. বাংলা অর্থ (\lat{pure})}

\textbf{\lat{“}তবে যারা তাওবা করেছে, সংশোধন করেছে এবং (সত্য) স্পষ্টভাবে বর্ণনা করেছে — তাদের তাওবা আমি গ্রহণ করব। আর আমিই তাওবা গ্রহণকারী, পরম দয়ালু।\lat{”}}

\subsubsection*{৪. শব্দের ব্যাখ্যা}

\begin{table}[h]\centering\small
\begin{tabular}{ll}
শব্দ & অর্থ \\
\hline
\textbf{আসলাহূ (\arab{أصلحوا})} & সংশোধন করেছে — তাওবার সাথে আমলের সংস্কার \\
\textbf{বাইয়্যানূ (\arab{بينوا})} & স্পষ্ট করেছে — গোপন করা সত্য প্রকাশ করা \\
\textbf{আতূবু 'আলাইহিম} & আমি তাদের তাওবা গ্রহণ করব \\
\hline
\end{tabular}
\end{table}

\subsubsection*{৫. সংক্ষিপ্ত তাফসীর}

\begin{itemize}
\item \textbf{ইবনে কাসীর:} এই আয়াতে তাওবার তিনটি স্তর দেখানো হয়েছে — (১) তাওবা করা, (২) আমল সংশোধন, (৩) গোপন সত্য প্রকাশ। অর্থাৎ — তাওবা মানে শুধু \lat{“}আফসোস\lat{”} নয়; সাথে সংশোধন ও সততা দরকার।
\item \textbf{সা'দী:} যারা আল্লাহর আয়াত গোপন করে মানুষকে ঠকিয়েছে, তাদের জন্যও দরজা খোলা — তাওবা করলে, সংশোধন হলে ও সত্য প্রকাশ করলে। এটি প্রমাণ — তাওবার দরজা কোনো শ্রেণির জন্য বন্ধ নয়।
\end{itemize}
\subsubsection*{৬. সংশ্লিষ্ট দলিল}

\begin{itemize}
\item \textbf{আয়াত (আল-ইমরান ৩:১৩৫-১৩৬):} পাপের পর ক্ষমা প্রার্থনাকারীদের বর্ণনা।
\end{itemize}
\medskip\hrule\medskip

\subsubsection*{আয়াত ১০ — সূরা আল-মা'ইদাহ (৫) — আয়াত ৩৯}

\subsubsection*{১. সূত্র কার্ড}

\begin{table}[h]\centering\small
\begin{tabular}{ll}
বিষয় & বিবরণ \\
\hline
\textbf{সূরা} & আল-মা'ইদাহ (৫) — মাদানী \\
\textbf{আয়াত} & ৩৯ \\
\textbf{মূল বিষয়} & অপরাধের পর তাওবা ও সংশোধন — আল্লাহ ক্ষমা করেন \\
\textbf{প্রসঙ্গ} & চোরের শাস্তির (হাত কাটা) পরের আয়াত — তাওবার পথ খোলা রাখা \\
\hline
\end{tabular}
\end{table}

\subsubsection*{২. মূল আরবি}

\begin{center}\arab{فَمَنْ تَابَ مِنْ بَعْدِ ظُلْمِهِ وَأَصْلَحَ فَإِنَّ اللَّهَ يَتُوبُ عَلَيْهِ ۗ إِنَّ اللَّهَ غَفُورٌ رَحِيمٌ}\end{center}
\textbf{উচ্চারণ:} ফামান তা-বা মিম বা'দি যুলমিহী ওয়া আসলাহা ফাইন্নাল্লা-হা ইয়াতূবু 'আলাইহি। ইন্নাল্লা-হা গাফূরুর রাহীম।

\subsubsection*{৩. বাংলা অর্থ (\lat{pure})}

\textbf{\lat{“}অতঃপর যে ব্যক্তি তার জুলুমের পর তাওবা করবে এবং সংশোধন করবে, নিশ্চয়ই আল্লাহ তার তাওবা কবুল করবেন। নিশ্চয়ই আল্লাহ ক্ষমাশীল, পরম দয়ালু।\lat{”}}

\subsubsection*{৪. শব্দের ব্যাখ্যা}

\begin{table}[h]\centering\small
\begin{tabular}{ll}
শব্দ & অর্থ \\
\hline
\textbf{যুলম (\arab{ظلم})} & অপরাধ/সীমালঙ্ঘন \\
\textbf{ইয়াতূবু 'আলাইহি} & তার তাওবা কবুল করেন — আল্লাহর পক্ষ থেকে তাওবার দিকে ফিরানো \\
\hline
\end{tabular}
\end{table}

\subsubsection*{৫. সংক্ষিপ্ত তাফসীর}

\begin{itemize}
\item \textbf{ইবনে কাসীর:} চুরির শাস্তি ঘোষণার পরপরই এই আয়াত — অর্থ: শাস্তি দুনিয়ার বিধান, কিন্তু তাওবার দরজা আল্লাহ খোলা রেখেছেন; যে তাওবা করে সে দুনিয়ার শাস্তি থেকে না পারলেও (শাস্তি কার্যকর হলে) আল্লাহর ক্ষমা পেতে পারে। এখানে আশার সবচেয়ে বড় দিক — শাস্তির ঘোষণাও তাওবার পথ বন্ধ করে না।
\item \textbf{সা'দী:} \lat{“}তাওবা করবে এবং সংশোধন করবে\lat{”} — দুই শর্ত; আবার শুধু \lat{“}তাওবা করল\lat{”} বললেও আল্লাহ ক্ষমা করেন — বড় অপরাধের পরেও।
\end{itemize}
\subsubsection*{৬. সংশ্লিষ্ট দলিল}

\begin{itemize}
\item \textbf{আয়াত (নিসা ৪:১১০):} \lat{“}যে মন্দ কাজ করবে বা নিজের উপর জুলুম করবে, তারপর ক্ষমা চাইবে — সে আল্লাহকে ক্ষমাশীল-দয়ালু পাবে।\lat{”}
\end{itemize}
\medskip\hrule\medskip

\subsubsection*{আয়াত ১১ — সূরা ত্বা-হা (২০) — আয়াত ৮২}

\subsubsection*{১. সূত্র কার্ড}

\begin{table}[h]\centering\small
\begin{tabular}{ll}
বিষয় & বিবরণ \\
\hline
\textbf{সূরা} & ত্বা-হা (২০) — মক্কী \\
\textbf{আয়াত} & ৮২ \\
\textbf{মূল বিষয়} & তাওবা, ঈমান ও সৎকাজ — আল্লাহর ক্ষমার শর্ত \\
\textbf{প্রসঙ্গ} & আদম (আঃ)-এর ঘটনার পর মানুষের প্রতি আল্লাহর ঘোষণা \\
\hline
\end{tabular}
\end{table}

\subsubsection*{২. মূল আরবি}

\begin{center}\arab{وَإِنِّي لَغَفَّارٌ لِمَنْ تَابَ وَآمَنَ وَعَمِلَ صَالِحًا ثُمَّ اهْتَدَىٰ}\end{center}
\textbf{উচ্চারণ:} ওয়া ইন্নী লাগাফফা-রুল লিমান তা-বা ওয়া আ-মানা ওয়া 'আমিলা সা-লিহান ছুম্মাহ তাদা-।

\subsubsection*{৩. বাংলা অর্থ (\lat{pure})}

\textbf{\lat{“}আর নিশ্চয়ই আমি অত্যন্ত ক্ষমাশীল তার জন্য, যে তাওবা করে, ঈমান আনে, সৎকাজ করে এবং এরপর সঠিক পথে থাকে।\lat{”}}

\subsubsection*{৪. শব্দের ব্যাখ্যা}

\begin{table}[h]\centering\small
\begin{tabular}{ll}
শব্দ & অর্থ \\
\hline
\textbf{লাগাফফার (\arab{لغفار})} & অত্যন্ত ক্ষমাশীল — বহুবার ক্ষমাকারী \\
\textbf{ছুম্মাহতাদা (\arab{ثم} \arab{اهتدى})} & এরপর হেদায়েতের উপর স্থির থাকা \\
\hline
\end{tabular}
\end{table}

\subsubsection*{৫. সংক্ষিপ্ত তাফসীর}

\begin{itemize}
\item \textbf{ইবনে কাসীর:} \lat{“}গাফফার\lat{”} — ক্ষমার সর্বোচ্চ মাত্রা; চারটি শর্তে ক্ষমা — তাওবা, ঈমান, সৎকাজ, স্থিরতা। অর্থ — তাওবা যাত্রার শুরু; এরপর চলতে থাকে ঈমান ও আমলের পথ।
\item \textbf{সা'দী:} তাওবার পর \lat{“}হেদায়েতের উপর স্থির থাকা\lat{”} — যেন তাওবার ফল ধরে রাখা হয়; ক্ষমার সাথে হেদায়েতের প্রতিশ্রুতি — আল্লাহ ক্ষমা করলেই হেদায়েতের পথ সুগম হয়।
\end{itemize}
\subsubsection*{৬. সংশ্লিষ্ট দলিল}

\begin{itemize}
\item \textbf{আয়াত (ফুরকান ২৫:৭০):} \lat{“}তাদের পাপগুলোকে আল্লাহ নেকিতে পরিবর্তন করে দেবেন।\lat{”}
\end{itemize}
\medskip\hrule\medskip

\subsubsection*{আয়াত ১২ — সূরা আন-নিসা (৪) — আয়াত ১১০}

\subsubsection*{১. সূত্র কার্ড}

\begin{table}[h]\centering\small
\begin{tabular}{ll}
বিষয় & বিবরণ \\
\hline
\textbf{সূরা} & আন-নিসা (৪) — মাদানী \\
\textbf{আয়াত} & ১১০ \\
\textbf{মূল বিষয়} & মন্দ কাজের পর ক্ষমা চাইলে আল্লাহ ক্ষমাশীল-দয়ালু পাবেই \\
\textbf{প্রসঙ্গ} & পাপের পর কী করা উচিত — তার সরাসরি নির্দেশনা \\
\hline
\end{tabular}
\end{table}

\subsubsection*{২. মূল আরবি}

\begin{center}\arab{وَمَنْ يَعْمَلْ سُوءًا أَوْ يَظْلِمْ نَفْسَهُ ثُمَّ يَسْتَغْفِرِ اللَّهَ يَجِدِ اللَّهَ غَفُورًا رَحِيمًا}\end{center}
\textbf{উচ্চারণ:} ওয়া মাই ইয়া'মাল সূআন আও ইয়াযলিম নাফসাহু ছুম্মা ইয়াসতাগফিরিল্লা-হা ইয়াজিদিল্লা-হা গাফূরার রাহীমা।

\subsubsection*{৩. বাংলা অর্থ (\lat{pure})}

\textbf{\lat{“}আর যে ব্যক্তি কোনো মন্দ কাজ করবে বা নিজের উপর জুলুম করবে, তারপর আল্লাহর কাছে ক্ষমা চাইবে — সে আল্লাহকে পাবে ক্ষমাশীল, পরম দয়ালু।\lat{”}}

\subsubsection*{৪. শব্দের ব্যাখ্যা}

\begin{table}[h]\centering\small
\begin{tabular}{ll}
শব্দ & অর্থ \\
\hline
\textbf{সূআন (\arab{سوءا})} & মন্দ কাজ \\
\textbf{ইয়াযলিমু নাফসাহু} & নিজের উপর জুলুম করা — নিজের ক্ষতি করা \\
\textbf{ইয়াসতাগফিরিল্লাহ} & আল্লাহর কাছে ক্ষমা প্রার্থনা করা \\
\hline
\end{tabular}
\end{table}

\subsubsection*{৫. সংক্ষিপ্ত তাফসীর}

\begin{itemize}
\item \textbf{ইবনে কাসীর:} এই আয়াত পাপের পরের করণীয় শেখায় — ভেঙে পড়া নয়, আল্লাহর কাছে ফিরে আসা; \lat{“}আল্লাহকে পাবে ক্ষমাশীল-দয়ালু\lat{”} — অর্থ ক্ষমা পেতে হলে তাকে খুঁজতে হবে, সে পাওয়া যায়।
\item \textbf{সা'দী:} পাপের পর দুটি পথ — হয় হতাশা (শয়তানের পথ), নয় ইস্তিগফার (আল্লাহর পথ); এই আয়াত দ্বিতীয় পথের নিশ্চয়তা দেয়।
\end{itemize}
\subsubsection*{৬. সংশ্লিষ্ট দলিল}

\begin{itemize}
\item \textbf{আয়াত (আল-ইমরান ৩:১৩৫):} \lat{“}…তারা আল্লাহকে স্মরণ করে এবং নিজেদের গুনাহের জন্য ক্ষমা চায়।\lat{”}
\item \textbf{হাদিস (বুখারি ৭৫০৭):} বারবার পাপ ও ক্ষমা প্রার্থনার ঘটনা।
\end{itemize}
\medskip\hrule\medskip

\subsection*{থিম ৪ — আল্লাহ ভালোবাসেন তাওবাকারীকে (\lat{Allah} \lat{Loves} \lat{the} \lat{Repentant})}

\subsubsection*{আয়াত ১৩ — সূরা আল-বাকারা (২) — আয়াত ২২২}

\subsubsection*{১. সূত্র কার্ড}

\begin{table}[h]\centering\small
\begin{tabular}{ll}
বিষয় & বিবরণ \\
\hline
\textbf{সূরা} & আল-বাকারা (২) — মাদানী \\
\textbf{আয়াত} & ২২২ \\
\textbf{মূল বিষয়} & আল্লাহ তাওবাকারীদের ভালোবাসেন — ভালোবাসার সরাসরি ঘোষণা \\
\textbf{প্রসঙ্গ} & হায়েযের বিধানের পর তাওবার কথা \\
\hline
\end{tabular}
\end{table}

\subsubsection*{২. মূল আরবি}

\begin{center}\arab{إِنَّ اللَّهَ يُحِبُّ التَّوَّابِينَ وَيُحِبُّ الْمُتَطَهِّرِينَ}\end{center}
\textbf{উচ্চারণ:} ইন্নাল্লা-হা ইউহিব্বুত তাওওয়াবীনা ওয়া ইউহিব্বুল মুতাত্বাহহিরীন।

\subsubsection*{৩. বাংলা অর্থ (\lat{pure})}

\textbf{\lat{“}নিশ্চয়ই আল্লাহ তাওবাকারীদের ভালোবাসেন এবং পবিত্রতা অবলম্বনকারীদের ভালোবাসেন।\lat{”}}

\subsubsection*{৪. শব্দের ব্যাখ্যা}

\begin{table}[h]\centering\small
\begin{tabular}{ll}
শব্দ & অর্থ \\
\hline
\textbf{ইউহিব্বু (\arab{يحب})} & ভালোবাসেন \\
\textbf{আত-তাওওয়াবীন (\arab{التوابين})} & তাওবাকারীরা — যারা বারবার তাওবা করে \\
\textbf{আল-মুতাত্বাহহিরীন} & পবিত্রতা অবলম্বনকারীরা \\
\hline
\end{tabular}
\end{table}

\subsubsection*{৫. সংক্ষিপ্ত তাফসীর}

\begin{itemize}
\item \textbf{ইবনে কাসীর:} \lat{“}তাওওয়াবীন\lat{”} — বহুবচন ও পুনরাবৃত্তির রূপ: যারা বারবার পাপ করে, বারবার তাওবা করে — আল্লাহ তাদের ভালোবাসেন। তাফসীরকারকদের মতে তাওবার এই স্তরই সর্বোচ্চ — পাপের পর লজ্জায় আল্লাহর কাছে ফিরে আসা।
\item \textbf{সা'দী:} ভালোবাসা শুধু \lat{“}ক্ষমা\lat{”} নয় — ভালোবাসা মানে সন্তুষ্টি ও বিশেষ রহমত; পাপী ব্যক্তি তাওবা করলে আল্লাহ শুধু ক্ষমা করেন না, ভালোবাসেন। এটাই পাপীর জন্য সবচেয়ে বড় আশা।
\end{itemize}
\subsubsection*{৬. সংশ্লিষ্ট দলিল}

\begin{itemize}
\item \textbf{হাদিস (মুসলিম ২৬৭৫):} \lat{“}আল্লাহ বান্দার তাওবায় এত খুশি হন…\lat{”}
\item \textbf{হাদিস (ইবনে মাজাহ ৪২৫১):} \lat{“}সেরা গুনাহকারী তাওবাকারী।\lat{”}
\end{itemize}
\medskip\hrule\medskip

\subsubsection*{আয়াত ১৪ — সূরা আল-ইমরান (৩) — আয়াত ১৩৫-১৩৬}

\subsubsection*{১. সূত্র কার্ড}

\begin{table}[h]\centering\small
\begin{tabular}{ll}
বিষয় & বিবরণ \\
\hline
\textbf{সূরা} & আল-ইমরান (৩) — মাদানী \\
\textbf{আয়াত} & ১৩৫-১৩৬ \\
\textbf{মূল বিষয়} & মুত্তাকীদের গুণ — পাপের পর সাথে সাথে ক্ষমা চাওয়া ও পাপে অটল না থাকা \\
\textbf{প্রসঙ্গ} & জান্নাতের জন্য প্রতিযোগীদের গুণাবলির বর্ণনা \\
\hline
\end{tabular}
\end{table}

\subsubsection*{২. মূল আরবি}

\begin{center}\arab{وَالَّذِينَ إِذَا فَعَلُوا فَاحِشَةً أَوْ ظَلَمُوا أَنْفُسَهُمْ ذَكَرُوا اللَّهَ فَاسْتَغْفَرُوا لِذُنُوبِهِمْ وَمَنْ يَغْفِرُ الذُّنُوبَ إِلَّا اللَّهُ وَلَمْ يُصِرُّوا عَلَىٰ مَا فَعَلُوا وَهُمْ يَعْلَمُونَ}\end{center}
\textbf{উচ্চারণ:} ওয়াল্লাযীনা ইযা- ফা'আলূ ফা-হিশাতান আও যালামূ আনফুসাহুম যাকারুল্লা-হা ফাসতাগফারূ লিজুনূবিহিম — ওয়া মাই ইয়াগফিরুয যুনূবা ইল্লাল্লা-হ — ওয়া লাম ইউসিররূ 'আলা- মা- ফা'আলূ ওয়া হুম ইয়া'লামূন।

\begin{center}\arab{أُولَٰئِكَ جَزَاؤُهُمْ مَغْفِرَةٌ مِنْ رَبِّهِمْ وَجَنَّاتٌ تَجْرِي مِنْ تَحْتِهَا الْأَنْهَارُ خَالِدِينَ فِيهَا ۚ وَنِعْمَ أَجْرُ الْعَامِلِينَ}\end{center}
\textbf{উচ্চারণ:} উলা-ইকা জাযাউহুম মাগফিরাতুম মির রব্বিহিম ওয়া জান্না-তুন তাজরী মিন তাহতিহাল আনহা-রু খা-লিদীনা ফীহা। ওয়া নি'মা আজরুল 'আ-মিলীন।

\subsubsection*{৩. বাংলা অর্থ (\lat{pure})}

\textbf{\lat{“}আর তারা — যখন কোনো অশ্লীল কাজ করে ফেলে বা নিজেদের উপর জুলুম করে, তখন আল্লাহকে স্মরণ করে এবং নিজেদের গুনাহের জন্য ক্ষমা চায়। আল্লাহ ছাড়া কে গুনাহ ক্ষমা করবে? আর তারা জেনে-বুঝে নিজেদের কৃতকর্মের উপর অটল থাকে না।\lat{”}} (৩:১৩৫)

\textbf{\lat{“}তাদের প্রতিদান — তাদের রবের পক্ষ থেকে ক্ষমা এবং জান্নাত, যার নিচে নদী প্রবাহিত; তারা সেখানে চিরকাল থাকবে। কতই না উত্তম আমলকারীদের পুরস্কার!\lat{”}} (৩:১৩৬)

\subsubsection*{৪. শব্দের ব্যাখ্যা}

\begin{table}[h]\centering\small
\begin{tabular}{ll}
শব্দ & অর্থ \\
\hline
\textbf{ফাহিশাহ (\arab{فاحشة})} & অশ্লীল/বড় পাপ \\
\textbf{লাম ইউসিররূ (\arab{لم} \arab{يصرّوا})} & অটল/স্থির থাকে না — পাপে ডুবে থাকে না \\
\textbf{ওয়া হুম ইয়া'লামূন} & জেনে-বুঝে — সচেতন অবস্থায় \\
\hline
\end{tabular}
\end{table}

\subsubsection*{৫. সংক্ষিপ্ত তাফসীর}

\begin{itemize}
\item \textbf{ইবনে কাসীর:} মুত্তাকীর একটি চিহ্ন — পাপ করলে সাথে সাথে আল্লাহকে স্মরণ ও ক্ষমা প্রার্থনা; আর \lat{“}পাপে অটল না থাকা\lat{”} তাওবার মূল — পাপ ছেড়ে দেওয়া। ইবনে কাসীর এখানে বলেছেন — পাপ করার সাথে সাথে তাওবা করা ফরজ; দেরি করলে পাপ জমে যায়।
\item \textbf{সা'দী:} \lat{“}আল্লাহ ছাড়া কে গুনাহ ক্ষমা করবে?\lat{”} — প্রশ্নসূচক বাক্যে ক্ষমার একমাত্র উৎস আল্লাহ — তাই অন্য কারো কাছে যাওয়ার দরকার নেই; সরাসরি তাঁর কাছে ফিরে আসতে হবে।
\item \textbf{মাআরিফুল কুরআন:} এই আয়াত মুত্তাকী ও পাপীকে এক পংক্তিতে এনেছে — কারণ মুত্তাকীও মানুষ, ভুল হয়; পার্থক্য শুধু — সে পাপের পর ফিরে আসে, পাপে অটল থাকে না।
\end{itemize}
\subsubsection*{৬. সংশ্লিষ্ট দলিল}

\begin{itemize}
\item \textbf{হাদিস (তিরমিযী ১৯৮৭):} \lat{“}মন্দ কাজের পর ভালো কাজ করো — তা তা মুছে দেবে।\lat{”}
\item \textbf{আয়াত (নিসা ৪:১১০):} ক্ষমা চাওয়ার নিশ্চয়তা।
\end{itemize}
\medskip\hrule\medskip

\subsection*{থিম ৫ — পাপ থেকে ফিরে: পরিবর্তন ও নতুন জীবন (\lat{New} \lat{Life} \lat{After} \lat{Sin})}

\subsubsection*{আয়াত ১৫ — সূরা আল-ফুরকান (২৫) — আয়াত ৭০-৭১ (সবচেয়ে শক্তিশালী আয়াতগুলোর একটি)}

\subsubsection*{১. সূত্র কার্ড}

\begin{table}[h]\centering\small
\begin{tabular}{ll}
বিষয় & বিবরণ \\
\hline
\textbf{সূরা} & আল-ফুরকান (২৫) — মক্কী \\
\textbf{আয়াত} & ৭০-৭১ \\
\textbf{মূল বিষয়} & তাওবার পর আল্লাহ পাপগুলোকে নেকিতে পরিবর্তন করে দেন — অতীত মুছে নতুন জীবন \\
\textbf{প্রসঙ্গ} & বড় পাপকারীদের (শিরক, হত্যা, যিনা) তাওবার ঘোষণা \\
\hline
\end{tabular}
\end{table}

\subsubsection*{২. মূল আরবি}

\begin{center}\arab{إِلَّا مَنْ تَابَ وَآمَنَ وَعَمِلَ عَمَلًا صَالِحًا فَأُولَٰئِكَ يُبَدِّلُ اللَّهُ سَيِّئَاتِهِمْ حَسَنَاتٍ ۗ وَكَانَ اللَّهُ غَفُورًا رَحِيمًا}\end{center}
\textbf{উচ্চারণ:} ইল্লা- মান তা-বা ওয়া আ-মানা ওয়া 'আমিলা 'আমালান সা-লিহান ফাউলা-ইকা ইউবাদ্দিলুল্লা-হু সাইয়ি-আতিহিম হাসানা-তিন। ওয়া কা-নাল্লা-হু গাফূরার রাহীমা।

\begin{center}\arab{وَمَنْ تَابَ وَعَمِلَ صَالِحًا فَإِنَّهُ يَتُوبُ إِلَى اللَّهِ مَتَابًا}\end{center}
\textbf{উচ্চারণ:} ওয়া মান তা-বা ওয়া 'আমিলা সা-লিহান ফাইন্নাহু ইয়াতূবু ইলাল্লা-হি মাতা-বা।

\subsubsection*{৩. বাংলা অর্থ (\lat{pure})}

\textbf{\lat{“}তবে যে তাওবা করেছে, ঈমান এনেছে এবং সৎকাজ করেছে — তাদের পাপগুলো আল্লাহ নেকিতে পরিবর্তন করে দেবেন। আল্লাহ ক্ষমাশীল, পরম দয়ালু।\lat{”}} (২৫:৭০)

\textbf{\lat{“}আর যে তাওবা করবে ও সৎকাজ করবে, সে তো আল্লাহর দিকে সঠিকভাবে ফিরে আসবে।\lat{”}} (২৫:৭১)

\subsubsection*{৪. শব্দের ব্যাখ্যা}

\begin{table}[h]\centering\small
\begin{tabular}{ll}
শব্দ & অর্থ \\
\hline
\textbf{ইউবাদ্দিলু (\arab{يبدّل})} & পরিবর্তন করে দেন — পাপের জায়গায় নেকি লেখা হয় \\
\textbf{সাইয়ি-আত (\arab{سيئات})} & খারাপ কাজ/পাপ \\
\textbf{মাতাবান (\arab{متابا})} & সঠিক/উত্তম প্রত্যাবর্তন — আল্লাহর দিকে ফেরার উত্তম পথ \\
\hline
\end{tabular}
\end{table}

\subsubsection*{৫. সংক্ষিপ্ত তাফসীর}

\begin{itemize}
\item \textbf{ইবনে কাসীর:} ২৫:৬৮-৬৯-এ শিরক, কুফরি, হত্যা ও যিনার কথা বলা হয়েছে; তারপর ৭০-এ তাওবার ঘোষণা — অর্থ: এই বড় বড় পাপের পরেও তাওবা কবুল। \lat{“}পাপকে নেকিতে পরিবর্তন\lat{”} — ইবনে আব্বাস ও অন্যরা ব্যাখ্যা করেছেন: কিয়ামতে তাদের পাপগুলোর স্থানে নেকি দেখানো হবে; ইবনে মাসউদের মতে — তাওবাকারীদের পাপ মুছে নেকি লেখা হয়। ইবনে কাসীর একে তাওবার সবচেয়ে বড় পুরস্কার বলেছেন।
\item \textbf{সা'দী:} এই আয়াত প্রমাণ করে — তাওবা শুধু পাপ মুছে দেয় না, বরং পাপের জায়গায় সৎকাজ বসিয়ে দেয়; অর্থাৎ তাওবাকারীর জীবনে বদল আসে — পুরনো পাপী চরিত্রের জায়গায় নতুন নেক চরিত্র।
\item \textbf{মাআরিফুল কুরআন:} মুশরিক, হত্যাকারী, যিনাকারী — সবাই তাওবায় ধরা পড়ে; ইসলামের পথে ফিরে এলে অতীতের ভয়ংকর অপরাধও আল্লাহর দয়ায় ধুয়ে যায়। এখান থেকেই পাপীর \lat{“}নতুন জীবন\lat{”}-এর কুরআনি ধারণা।
\end{itemize}
\subsubsection*{৬. সংশ্লিষ্ট দলিল}

\begin{itemize}
\item \textbf{হাদিস (মুসলিম ২৭৬৬):} ৯৯+১ জন হত্যাকারীর তাওবা — এই আয়াতেরই জীবন্ত উদাহরণ।
\item \textbf{আয়াত (যুমার ৩৯:৫৩):} ক্ষমার সর্বোচ্চ ঘোষণা।
\end{itemize}
\medskip\hrule\medskip

\subsubsection*{আয়াত ১৬ — সূরা আন-নাহল (১৬) — আয়াত ১১৯}

\subsubsection*{১. সূত্র কার্ড}

\begin{table}[h]\centering\small
\begin{tabular}{ll}
বিষয় & বিবরণ \\
\hline
\textbf{সূরা} & আন-নাহল (১৬) — মক্কী \\
\textbf{আয়াত} & ১১৯ \\
\textbf{মূল বিষয়} & অজ্ঞতাবশত পাপের পর তাওবা ও সংশোধন — ক্ষমার নিশ্চয়তা \\
\textbf{প্রসঙ্গ} & আগের আয়াতের (অজ্ঞতায় পাপকারীদের) পরিপূরক ঘোষণা \\
\hline
\end{tabular}
\end{table}

\subsubsection*{২. মূল আরবি}

\begin{center}\arab{ثُمَّ إِنَّ رَبَّكَ لِلَّذِينَ عَمِلُوا السُّوءَ بِجَهَالَةٍ ثُمَّ تَابُوا مِنْ بَعْدِ ذَٰلِكَ وَأَصْلَحُوا إِنَّ رَبَّكَ مِنْ بَعْدِهَا لَغَفُورٌ رَحِيمٌ}\end{center}
\textbf{উচ্চারণ:} ছুম্মা ইন্না রব্বাকা লিল্লাযীনা 'আমিলুস সূআ বিজাহা-লাতিন ছুম্মা তা-বূ মিম বা'দি যা-লিকা ওয়া আসলাহূ, ইন্না রব্বাকা মিম বা'দিহা- লাগাফূরুর রাহীম।

\subsubsection*{৩. বাংলা অর্থ (\lat{pure})}

\textbf{\lat{“}অতঃপর নিশ্চয়ই তোমার রব — যারা অজ্ঞতাবশত মন্দ কাজ করেছে, এরপর তাওবা করেছে এবং সংশোধন করেছে — নিশ্চয়ই তাদের জন্য এরপর তোমার রব ক্ষমাশীল, পরম দয়ালু।\lat{”}}

\subsubsection*{৪. শব্দের ব্যাখ্যা}

\begin{table}[h]\centering\small
\begin{tabular}{ll}
শব্দ & অর্থ \\
\hline
\textbf{বিজাহালাতিন (\arab{بجهالة})} & অজ্ঞতাবশত — তাফসীরে: যে কোনো পাপই আল্লাহ সম্পর্কে অজ্ঞতা থেকে হয় \\
\textbf{আসলাহূ (\arab{أصلحوا})} & সংশোধন করেছে \\
\hline
\end{tabular}
\end{table}

\subsubsection*{৫. সংক্ষিপ্ত তাফসীর}

\begin{itemize}
\item \textbf{ইবনে কাসীর:} কোনো পাপই জ্ঞান-সচেতনতা দিয়ে হয় না — সব পাপই \lat{“}জাহালাত\lat{”} (অজ্ঞতা) থেকে; তাই এই শব্দে সব পাপী ধরা পড়ে। তাওবা ও সংশোধন করলেই ক্ষমা।
\item \textbf{সা'দী:} দুই শর্ত — তাওবা (অতীত ত্যাগ) ও সংশোধন (ভবিষ্যত সংস্কার); দুই-ই হলে \lat{“}ক্ষমাশীল-দয়ালু\lat{”} পাওয়া নিশ্চিত।
\end{itemize}
\subsubsection*{৬. সংশ্লিষ্ট দলিল}

\begin{itemize}
\item \textbf{আয়াত (আন'আম ৬:৫৪):} \lat{“}যে অজ্ঞতাবশত মন্দ করল, তারপর তাওবা ও সংশোধন করল…\lat{”}
\item \textbf{আয়াত (ফুরকান ২৫:৭০):} পাপের স্থানে নেকি।
\end{itemize}
\medskip\hrule\medskip

\subsubsection*{আয়াত ১৭ — সূরা আল-আ'রাফ (৭) — আয়াত ১৫৩}

\subsubsection*{১. সূত্র কার্ড}

\begin{table}[h]\centering\small
\begin{tabular}{ll}
বিষয় & বিবরণ \\
\hline
\textbf{সূরা} & আল-আ'রাফ (৭) — মক্কী \\
\textbf{আয়াত} & ১৫৩ \\
\textbf{মূল বিষয়} & গো-বাছুর পূজার (বড় পাপ) পরেও তাওবার দরজা \\
\textbf{প্রসঙ্গ} & বনি ইসরাঈলের গো-বাছুর পূজার পর তাওবার ঘোষণা \\
\hline
\end{tabular}
\end{table}

\subsubsection*{২. মূল আরবি}

\begin{center}\arab{وَالَّذِينَ عَمِلُوا السَّيِّئَاتِ ثُمَّ تَابُوا مِنْ بَعْدِهَا وَآمَنُوا إِنَّ رَبَّكَ مِنْ بَعْدِهَا لَغَفُورٌ رَحِيمٌ}\end{center}
\textbf{উচ্চারণ:} ওয়াল্লাযীনা 'আমিলুস সাইয়ি-আতি ছুম্মা তা-বূ মিম বা'দিহা- ওয়া আ-মানূ, ইন্না রব্বাকা মিম বা'দিহা- লাগাফূরুর রাহীম।

\subsubsection*{৩. বাংলা অর্থ (\lat{pure})}

\textbf{\lat{“}আর যারা মন্দ কাজ করেছে, এরপর তাওবা করেছে ও ঈমান এনেছে — নিশ্চয়ই এরপর তোমার রব ক্ষমাশীল, পরম দয়ালু।\lat{”}}

\subsubsection*{৪. সংক্ষিপ্ত তাফসীর}

\begin{itemize}
\item \textbf{ইবনে কাসীর:} এই আয়াত গো-বাছুর পূজার পরের — অর্থ: শিরকের মতো পাপও তাওবার মাধ্যমে ক্ষমা করা হয়েছে; তবে তাওবা হতে হবে মৃত্যুর পূর্বে ও ঈমানের সাথে।
\item \textbf{সা'দী:} তাওবা + ঈমান = ক্ষমা; এখানে \lat{“}ঈমান\lat{”} যোগ হয়েছে — তাওবার সাথে আকীদার সংশোধন।
\end{itemize}
\subsubsection*{৫. সংশ্লিষ্ট দলিল}

\begin{itemize}
\item \textbf{আয়াত (যুমার ৩৯:৫৩):} সব গুনাহের ক্ষমার ঘোষণা।
\end{itemize}
\medskip\hrule\medskip

\subsubsection*{আয়াত ১৮ — সূরা আল-ইসরা (১৭) — আয়াত ২৫}

\subsubsection*{১. সূত্র কার্ড}

\begin{table}[h]\centering\small
\begin{tabular}{ll}
বিষয় & বিবরণ \\
\hline
\textbf{সূরা} & আল-ইসরা (১৭) — মক্কী \\
\textbf{আয়াত} & ২৫ \\
\textbf{মূল বিষয়} & আল্লাহর কাছে বারবার ফিরে আসা (আওওয়াব) — ক্ষমার প্রতিশ্রুতি \\
\textbf{প্রসঙ্গ} & পিতামাতার প্রতি দয়ার নির্দেশনার পর অন্তরের কথা \\
\hline
\end{tabular}
\end{table}

\subsubsection*{২. মূল আরবি}

\begin{center}\arab{رَبُّكُمْ أَعْلَمُ بِمَا فِي نُفُوسِكُمْ ۚ إِنْ تَكُونُوا صَالِحِينَ فَإِنَّهُ كَانَ لِلْأَوَّابِينَ غَفُورًا}\end{center}
\textbf{উচ্চারণ:} রব্বুকুম আ'লামু বিমা- ফী নুফূসিকুম। ইন তাকূনূ সা-লিহীনা ফাইন্নাহু কা-না লিলআওওয়াবীনা গাফূরা।

\subsubsection*{৩. বাংলা অর্থ (\lat{pure})}

\textbf{\lat{“}তোমাদের রব তোমাদের অন্তরে যা আছে তা ভালো জানেন। যদি তোমরা সৎ হও, তবে নিশ্চয়ই তিনি বারবার তাঁর দিকে প্রত্যাবর্তনকারীদের (আওওয়াবীন) জন্য ক্ষমাশীল।\lat{”}}

\subsubsection*{৪. শব্দের ব্যাখ্যা}

\begin{table}[h]\centering\small
\begin{tabular}{ll}
শব্দ & অর্থ \\
\hline
\textbf{আল-আওওয়াব (\arab{الأواب})} & যে বারবার আল্লাহর দিকে ফিরে আসে — পাপের পর তাওবা, আবার পাপের পর আবার তাওবা \\
\hline
\end{tabular}
\end{table}

\subsubsection*{৫. সংক্ষিপ্ত তাফসীর}

\begin{itemize}
\item \textbf{ইবনে কাসীর:} \lat{“}আওওয়াব\lat{”} এমন ব্যক্তি — যে গুনাহের পর তাওবা করে, আবার ভুল করলে আবার তাওবা করে; ইবনে আব্বাস বলেন — আওওয়াব মানে জিকিরকারী ও তাওবাকারী।
\item \textbf{সা'দী:} অন্তরের অবস্থা আল্লাহ জানেন — তাই শুধু বাহ্যিক তাওবা নয়; অন্তর থেকে ফেরা দরকার। আর যারা বারবার ফেরে, তাদের জন্য ক্ষমার বিশেষ প্রতিশ্রুতি।
\end{itemize}
\subsubsection*{৬. সংশ্লিষ্ট দলিল}

\begin{itemize}
\item \textbf{আয়াত (বাকারা ২:২২২):} \lat{“}আল্লাহ তাওবাকারীদের ভালোবাসেন।\lat{”}
\item \textbf{হাদিস (মুসলিম ২৭৫৯):} বারবার পাপ-তাওবার অধ্যায়ের হাদিস।
\end{itemize}
\medskip\hrule\medskip

\subsubsection*{আয়াত ১৯ — সূরা আন-নাজম (৫৩) — আয়াত ৩২}

\subsubsection*{১. সূত্র কার্ড}

\begin{table}[h]\centering\small
\begin{tabular}{ll}
বিষয় & বিবরণ \\
\hline
\textbf{সূরা} & আন-নাজম (৫৩) — মক্কী \\
\textbf{আয়াত} & ৩২ \\
\textbf{মূল বিষয়} & বড় পাপ থেকে বিরত থাকার পর ছোট ভুলের ক্ষমা — রবের ক্ষমা প্রশস্ত \\
\textbf{প্রসঙ্গ} & মুত্তাকীদের প্রশংসার প্রসঙ্গে \\
\hline
\end{tabular}
\end{table}

\subsubsection*{২. মূল আরবি}

\begin{center}\arab{الَّذِينَ يَجْتَنِبُونَ كَبَائِرَ الْإِثْمِ وَالْفَوَاحِشَ إِلَّا اللَّمَمَ ۚ إِنَّ رَبَّكَ وَاسِعُ الْمَغْفِرَةِ}\end{center}
\textbf{উচ্চারণ:} আল্লাযীনা ইয়াজতানিবূনা কাবা-ইরাল ইসমি ওয়াল ফাওয়া-হিশা ইল্লাল লামাম। ইন্না রব্বাকা ওয়া-সিউল মাগফিরাহ।

\subsubsection*{৩. বাংলা অর্থ (\lat{pure})}

\textbf{\lat{“}যারা বড় পাপ ও অশ্লীল কাজ থেকে বিরত থাকে — তবে ছোটখাটো (ভুল) ছাড়া। নিশ্চয়ই তোমার রবের ক্ষমা প্রশস্ত।\lat{”}}

\subsubsection*{৪. শব্দের ব্যাখ্যা}

\begin{table}[h]\centering\small
\begin{tabular}{ll}
শব্দ & অর্থ \\
\hline
\textbf{কাবা-ইরুল ইসম (\arab{كبائر} \arab{الإثم})} & বড় পাপ \\
\textbf{আল-লামাম (\arab{اللمم})} & ছোটখাটো পাপ/ভুল — বারবার নয় এমন সামান্য অপরাধ \\
\textbf{ওয়াসিউল মাগফিরাহ (\arab{واسع} \arab{المغفرة})} & প্রশস্ত ক্ষমা — ব্যাপক ক্ষমার মালিক \\
\hline
\end{tabular}
\end{table}

\subsubsection*{৫. সংক্ষিপ্ত তাফসীর}

\begin{itemize}
\item \textbf{ইবনে কাসীর:} মুত্তাকীও ছোটখাটো ভুল করেন — কিন্তু বড় পাপ ও অশ্লীলতা থেকে বিরত থাকেন; আর আল্লাহর ক্ষমা প্রশস্ত — ছোট ভুল ক্ষমার আওতায়। এই আয়াত থেকে বোঝা যায় — পূর্ণতা না থাকলেও আশা আছে; আল্লাহ ছোট ভুলকে ক্ষমা করেন, যদি বড় পাপে ডুবে না থাকে।
\item \textbf{সা'দী:} ক্ষমার প্রশস্ততা বান্দার পাপের আকারের চেয়ে বড় — তাই কেউ যেন ক্ষমা নিয়ে হতাশ না হয়।
\end{itemize}
\subsubsection*{৬. সংশ্লিষ্ট দলিল}

\begin{itemize}
\item \textbf{আয়াত (নিসা ৪:৪৮):} শিরক ছাড়া অন্য সব পাপ ক্ষমার আওতায় (আল্লাহ চাইলে)।
\end{itemize}
\medskip\hrule\medskip

\subsubsection*{আয়াত ২০ — সূরা হুদ (১১) — আয়াত ১১৪}

\subsubsection*{১. সূত্র কার্ড}

\begin{table}[h]\centering\small
\begin{tabular}{ll}
বিষয় & বিবরণ \\
\hline
\textbf{সূরা} & হুদ (১১) — মক্কী \\
\textbf{আয়াত} & ১১৪ \\
\textbf{মূল বিষয়} & ভালো কাজ মন্দ কাজ মুছে দেয় — তাওবার পরের আমলের পথ \\
\textbf{প্রসঙ্গ} & নামাজের গুরুত্বের প্রসঙ্গে \\
\hline
\end{tabular}
\end{table}

\subsubsection*{২. মূল আরবি}

\begin{center}\arab{إِنَّ الْحَسَنَاتِ يُذْهِبْنَ السَّيِّئَاتِ ۚ ذَٰلِكَ ذِكْرَىٰ لِلذَّاكِرِينَ}\end{center}
\textbf{উচ্চারণ:} ইন্নাল হাসানা-তি ইউযহিবনাস সাইয়ি-আ-তি। যা-লিকা যিকরা- লিয যা-কিরীন।

\subsubsection*{৩. বাংলা অর্থ (\lat{pure})}

\textbf{\lat{“}নিশ্চয়ই ভালো কাজগুলো মন্দ কাজগুলোকে মুছে দেয়। এটা স্মরণকারীদের জন্য উপদেশ।\lat{”}}

\subsubsection*{৪. শব্দের ব্যাখ্যা}

\begin{table}[h]\centering\small
\begin{tabular}{ll}
শব্দ & অর্থ \\
\hline
\textbf{ইউযহিবনা (\arab{يذهبن})} & মুছে দেয়/দূর করে দেয় \\
\textbf{আল-হাসানাত} & ভালো কাজ — বিশেষত নামাজ (প্রসঙ্গ অনুযায়ী) \\
\hline
\end{tabular}
\end{table}

\subsubsection*{৫. সংক্ষিপ্ত তাফসীর}

\begin{itemize}
\item \textbf{ইবনে কাসীর:} আয়াতের ব্যাখ্যায় নবীজি (সা.) বলেছেন — পাঁচ ওয়াক্ত নামাজ, জুমা থেকে জুমা, রমজান থেকে রমজান — মাঝের পাপ মুছে দেয়, বড় পাপ ছাড়া। এছাড়া হাদিসে এসেছে — \lat{“}মন্দ কাজের পর ভালো কাজ করো\lat{”} (তিরমিযী ১৯৮৭)। অর্থ — তাওবার সাথে ভালো আমল পাপের চিহ্ন মুছে দেয়।
\item \textbf{সা'দী:} ভালো কাজ শুধু পাপ মুছে দেয় না — অন্তরকে পরিষ্কার করে, নতুন জীবনের পথ খুলে দেয়।
\end{itemize}
\subsubsection*{৬. সংশ্লিষ্ট দলিল}

\begin{itemize}
\item \textbf{হাদিস (তিরমিযী ১৯৮৭):} \lat{“}তুমি যেখানে দাঁড়িয়ে আছ, মন্দ কাজের পর ভালো কাজ করো — তা মুছে দেবে।\lat{”}
\item \textbf{আয়াত (ফুরকান ২৫:৭০):} পাপের স্থানে নেকি।
\end{itemize}
\medskip\hrule\medskip

\subsubsection*{আয়াত ২১ — সূরা আত-তাওবা (৯) — আয়াত ১০২}

\subsubsection*{১. সূত্র কার্ড}

\begin{table}[h]\centering\small
\begin{tabular}{ll}
বিষয় & বিবরণ \\
\hline
\textbf{সূরা} & আত-তাওবা (৯) — মাদানী \\
\textbf{আয়াত} & ১০২ \\
\textbf{মূল বিষয়} & যারা নিজেদের পাপ স্বীকার করেছে — ভালো-মন্দ মিশ্রিত আমল — আল্লাহ তাদের প্রতি ফিরবেন \\
\textbf{প্রসঙ্গ} & তাবুকের যুদ্ধে পিছিয়ে পড়া কিছু সাহাবি — তাওবার অপেক্ষায় \\
\hline
\end{tabular}
\end{table}

\subsubsection*{২. মূল আরবি}

\begin{center}\arab{وَآخَرُونَ اعْتَرَفُوا بِذُنُوبِهِمْ خَلَطُوا عَمَلًا صَالِحًا وَآخَرَ سَيِّئًا عَسَى اللَّهُ أَنْ يَتُوبَ عَلَيْهِمْ ۚ إِنَّ اللَّهَ غَفُورٌ رَحِيمٌ}\end{center}
\textbf{উচ্চারণ:} ওয়া আ-খারূনা' তারাফূ বিজুনূবিহিম খালাতূ 'আমালান সা-লিহান ওয়া আ-খারা সাইয়ি-আন, 'আসাল্লা-হু আই ইয়াতূবা 'আলাইহিম। ইন্নাল্লা-হা গাফূরুর রাহীম।

\subsubsection*{৩. বাংলা অর্থ (\lat{pure})}

\textbf{\lat{“}আর অন্যেরা আছে যারা নিজেদের পাপ স্বীকার করেছে — তারা ভালো ও মন্দ — মিশ্রিত আমল করেছে। আশা করা যায়, আল্লাহ তাদের প্রতি ক্ষমার দৃষ্টি দেবেন। নিশ্চয়ই আল্লাহ ক্ষমাশীল, পরম দয়ালু।\lat{”}}

\subsubsection*{৪. শব্দের ব্যাখ্যা}

\begin{table}[h]\centering\small
\begin{tabular}{ll}
শব্দ & অর্থ \\
\hline
\textbf{ই'তারাফূ (\arab{اعترفوا})} & স্বীকার করেছে — পাপ স্বীকার করা তাওবার প্রথম ধাপ \\
\textbf{খালাতূ (\arab{خلطوا})} & মিশ্রিত করেছে — ভালো-মন্দ দুই-ই আছে \\
\textbf{আসাল্লাহু আই ইয়াতূবা 'আলাইহিম} & আশা করা যায় আল্লাহ তাদের প্রতি ফিরবেন \\
\hline
\end{tabular}
\end{table}

\subsubsection*{৫. সংক্ষিপ্ত তাফসীর}

\begin{itemize}
\item \textbf{ইবনে কাসীর:} এই আয়াত তাবুকের ঘটনায় পিছিয়ে পড়া সাহাবিদের ব্যাপারে — যারা পাপ স্বীকার করেছিল; আল্লাহ বললেন \lat{“}আশা করা যায় ফিরব\lat{”} — এরপর ৯:১১৭-১১৮-এ তাদের তাওবা কবুলের ঘোষণা। শিক্ষা — পাপ স্বীকার করাই আল্লাহর রহমত টানার প্রথম ধাপ; গোপন রাখা বা অজুহাত দেওয়া নয়।
\item \textbf{সা'দী:} ভালো-মন্দ মিশ্রিত জীবনও আশার বাইরে নয় — যদি অন্তরে পাপের বোঝা ও আল্লাহর দিকে ফেরার ইচ্ছা থাকে।
\end{itemize}
\subsubsection*{৬. সংশ্লিষ্ট দলিল}

\begin{itemize}
\item \textbf{আয়াত (তাওবা ৯:১১৮):} তিন সাহাবির তাওবা কবুলের পূর্ণ ঘোষণা।
\item \textbf{ঘটনা:} কাব ইবনে মালিক (রাঃ) — সাহাবিদের ঘটনা ফাইল ১।
\end{itemize}
\medskip\hrule\medskip

\subsection*{থিম ৬ — হেদায়েত, সাহায্য ও দোয়ার উত্তর (\lat{Guidance} \lat{\&} \lat{Help})}

\subsubsection*{আয়াত ২২ — সূরা আল-আনকাবুত (২৯) — আয়াত ৬৯}

\subsubsection*{১. সূত্র কার্ড}

\begin{table}[h]\centering\small
\begin{tabular}{ll}
বিষয় & বিবরণ \\
\hline
\textbf{সূরা} & আল-আনকাবুত (২৯) — মক্কী \\
\textbf{আয়াত} & ৬৯ \\
\textbf{মূল বিষয়} & যারা আল্লাহর পথে সংগ্রাম করে, আল্লাহ তাদের পথ দেখান — সাহায্যের প্রতিশ্রুতি \\
\textbf{প্রসঙ্গ} & ঈমান ও সংগ্রামের ফল \\
\hline
\end{tabular}
\end{table}

\subsubsection*{২. মূল আরবি}

\begin{center}\arab{وَالَّذِينَ جَاهَدُوا فِينَا لَنَهْدِيَنَّهُمْ سُبُلَنَا ۚ وَإِنَّ اللَّهَ لَمَعَ الْمُحْسِنِينَ}\end{center}
\textbf{উচ্চারণ:} ওয়াল্লাযীনা জা-হাদূ ফীনা- লানাহদিয়ান্নাহুম সুলানা-। ওয়া ইন্নাল্লা-হা লামা'আল মুহসিনীন।

\subsubsection*{৩. বাংলা অর্থ (\lat{pure})}

\textbf{\lat{“}আর যারা আমার পথে সংগ্রাম করে, আমি অবশ্যই তাদেরকে আমার পথসমূহ দেখিয়ে দেব। আর নিশ্চয়ই আল্লাহ সৎকর্মপরায়ণদের সাথে আছেন।\lat{”}}

\subsubsection*{৪. শব্দের ব্যাখ্যা}

\begin{table}[h]\centering\small
\begin{tabular}{ll}
শব্দ & অর্থ \\
\hline
\textbf{জা-হাদূ (\arab{جاهدوا})} & সংগ্রাম করেছে — নফসের সাথে, পাপের সাথে লড়াই \\
\textbf{লানাহদিয়ান্নাহুম সুলানা (\arab{لنهدينهم} \arab{سبلنا})} & অবশ্যই তাদের আমাদের পথ দেখাব \\
\hline
\end{tabular}
\end{table}

\subsubsection*{৫. সংক্ষিপ্ত তাফসীর}

\begin{itemize}
\item \textbf{ইবনে কাসীর:} যে ব্যক্তি আল্লাহর সন্তুষ্টির জন্য সংগ্রাম করে — পাপ ত্যাগ, তাওবা, আমলের পথে — আল্লাহ তাকে হেদায়েতের নতুন দরজা খুলে দেন। তাফসীরকারকগণ বলেন — প্রতিশ্রুতি \lat{“}অবশ্যই\lat{”} (\arab{لن}) দিয়ে — কোনো শর্ত ছাড়া নিশ্চিত।
\item \textbf{সা'দী:} হেদায়েতই সবচেয়ে বড় সাহায্য — পাপী ব্যক্তি যখন তাওবার জন্য সত্যিই সংগ্রাম করে (মানে চেষ্টা), আল্লাহ তাকে পথ দেখান; এই আয়াত পাপীর জন্য আশার বড় দলিল — ফিরে আসার চেষ্টাই আল্লাহর সাহায্য টানে।
\end{itemize}
\subsubsection*{৬. সংশ্লিষ্ট দলিল}

\begin{itemize}
\item \textbf{আয়াত (বাকারা ২:১৮৬):} \lat{“}আমি নিকটে — ডাকলে সাড়া দিই।\lat{”}
\item \textbf{হাদিস (মুসলিম ২৬৭৫):} \lat{“}আমার দিকে এক বিঘত এগোলো — আমি এক হাত এগিয়ে আসি…\lat{”}
\end{itemize}
\medskip\hrule\medskip

\subsubsection*{আয়াত ২৩ — সূরা আল-বাকারা (২) — আয়াত ১৮৬}

\subsubsection*{১. সূত্র কার্ড}

\begin{table}[h]\centering\small
\begin{tabular}{ll}
বিষয় & বিবরণ \\
\hline
\textbf{সূরা} & আল-বাকারা (২) — মাদানী \\
\textbf{আয়াত} & ১৮৬ \\
\textbf{মূল বিষয়} & আল্লাহ নিকটে — দোয়ার উত্তর দেওয়ার ঘোষণা \\
\textbf{প্রসঙ্গ} & সাহাবিরা জিজ্ঞেস করেছিলেন — আল্লাহ কোথায়, কীভাবে ডাকব? \\
\hline
\end{tabular}
\end{table}

\subsubsection*{২. মূল আরবি}

\begin{center}\arab{وَإِذَا سَأَلَكَ عِبَادِي عَنِّي فَإِنِّي قَرِيبٌ ۖ أُجِيبُ دَعْوَةَ الدَّاعِ إِذَا دَعَانِ ۖ فَلْيَسْتَجِيبُوا لِي وَلْيُؤْمِنُوا بِي لَعَلَّهُمْ يَرْشُدُونَ}\end{center}
\textbf{উচ্চারণ:} ওয়া ইযা- সাআলাকা 'ইবা-দী 'আন্নী ফাইন্নী কারীব। উজীবু দা'ওয়াতাদ দা-'ই ইযা- দা'আ-নি। ফালইয়াসতাজীবূ লী ওয়া লইউ'মিনূ বী লা'আল্লাহুম ইয়ারশুদূন।

\subsubsection*{৩. বাংলা অর্থ (\lat{pure})}

\textbf{\lat{“}আর যখন আমার বান্দারা তোমাকে আমার সম্পর্কে জিজ্ঞেস করে — আমি তো নিকটেই আছি; ডাক্তারের ডাকে সাড়া দিই, যখন সে আমাকে ডাকে। সুতরাং তারা যেন আমার ডাকে সাড়া দেয় এবং আমার প্রতি ঈমান রাখে — যাতে তারা সঠিক পথ পায়।\lat{”}}

\subsubsection*{৪. শব্দের ব্যাখ্যা}

\begin{table}[h]\centering\small
\begin{tabular}{ll}
শব্দ & অর্থ \\
\hline
\textbf{ফাইন্নী কারীব (\arab{فإني} \arab{قريب})} & আমি তো নিকটে — দূরত্ব নেই \\
\textbf{উজীবু দা'ওয়াতাদ দা'ই} & ডাকের জবাব দিই \\
\hline
\end{tabular}
\end{table}

\subsubsection*{৫. সংক্ষিপ্ত তাফসীর}

\begin{itemize}
\item \textbf{ইবনে কাসীর:} এই আয়াত নাজিলের প্রেক্ষাপট — সাহাবিরা জিজ্ঞেস করলেন, আমাদের রব কি নিকটে (তাহলে ফিসফিস করে বলব) নাকি দূরে (তাহলে চিৎকার করে ডাকব)? আল্লাহ বললেন — আমি নিকটে; ডাকার পদ্ধতি — নীরবে, বিনয়ে, অন্তর থেকে। পাপী ব্যক্তির জন্য — যত বড় পাপই হোক, ডাকার দূরত্ব নেই; লজ্জায় ডাকা বন্ধ করার কোনো কারণ নেই।
\item \textbf{সা'দী:} আয়াতের শেষাংশ — \lat{“}তারা যেন আমার ডাকে সাড়া দেয়\lat{”} — অর্থ: তাওবা ও ইবাদতের ডাকে; আল্লাহ ডাকেন তাওবার দিকে, বান্দা সাড়া দিলে উত্তরে আল্লাহ তার দোয়ার জবাব দেন।
\end{itemize}
\subsubsection*{৬. সংশ্লিষ্ট দলিল}

\begin{itemize}
\item \textbf{আয়াত (গাফির ৪০:৬০):} \lat{“}তোমরা আমাকে ডাকো, আমি সাড়া দেব।\lat{”}
\item \textbf{আয়াত (বাকারা ২:১৫২):} \lat{“}আমাকে স্মরণ করো, আমি তোমাদের স্মরণ করব।\lat{”}
\end{itemize}
\medskip\hrule\medskip

\subsubsection*{আয়াত ২৪ — সূরা আল-মু'মিন (৪০) — আয়াত ৬০}

\subsubsection*{১. সূত্র কার্ড}

\begin{table}[h]\centering\small
\begin{tabular}{ll}
বিষয় & বিবরণ \\
\hline
\textbf{সূরা} & আল-মু'মিন (গাফির) (৪০) — মক্কী \\
\textbf{আয়াত} & ৬০ \\
\textbf{মূল বিষয়} & দোয়ার প্রতিশ্রুতি — ডাকলে উত্তর \\
\textbf{প্রসঙ্গ} & ঈমানদারদের দোয়ার নির্দেশ \\
\hline
\end{tabular}
\end{table}

\subsubsection*{২. মূল আরবি}

\begin{center}\arab{وَقَالَ رَبُّكُمُ ادْعُونِي أَسْتَجِبْ لَكُمْ ۚ إِنَّ الَّذِينَ يَسْتَكْبِرُونَ عَنْ عِبَادَتِي سَيَدْخُلُونَ جَهَنَّمَ دَاخِرِينَ}\end{center}
\textbf{উচ্চারণ:} ওয়া কা-লা রব্বুকুমুদ 'ঊনী আসতাজিব লাকুম। ইন্নাল্লাযীনা ইয়াসতাকবিরূনা 'আন 'ইবা-দতী সাইয়াদখুলূনা জাহান্নামা দা-খিরীন।

\subsubsection*{৩. বাংলা অর্থ (\lat{pure})}

\textbf{\lat{“}আর তোমাদের রব বলেছেন: তোমরা আমাকে ডাকো, আমি তোমাদের ডাকে সাড়া দেব। নিশ্চয়ই যারা আমার ইবাদত থেকে অহংকার করে (বিমুখ থাকে), তারা লাঞ্ছিত হয়ে জাহান্নামে প্রবেশ করবে।\lat{”}}

\subsubsection*{৪. সংক্ষিপ্ত তাফসীর}

\begin{itemize}
\item \textbf{ইবনে কাসীর:} দোয়া ইবাদতের সারমর্ম; আল্লাহ নিজেই ডাকার আমন্ত্রণ দিয়েছেন — প্রতিশ্রুতি দিয়েছেন সাড়া দেওয়ার। পাপী ব্যক্তি দোয়া করলে আল্লাহ ফিরিয়ে দেন না — প্রতিশ্রুতি অনুযায়ী জবাব আসে (তাঁর হিকমত অনুযায়ী)।
\item \textbf{সা'দী:} দোয়ার মধ্যে সবচেয়ে বড় দোয়া — ক্ষমা ও হেদায়েতের দোয়া; পাপী যখন এ দোয়ায় ফিরে আসে, সে ইবাদতের মূলেই ফিরে আসে।
\end{itemize}
\subsubsection*{৫. সংশ্লিষ্ট দলিল}

\begin{itemize}
\item \textbf{আয়াত (বাকারা ২:১৮৬):} \lat{“}আমি নিকটে, ডাকলে সাড়া দিই।\lat{”}
\item \textbf{হাদিস (তিরমিযী ৩৫০৫):} ইউনুস (আঃ)-এর দোয়া — \lat{“}কোনো মুসলিম তা দিয়ে দোয়া করলে আল্লাহ সাড়া দেন।\lat{”}
\end{itemize}
\medskip\hrule\medskip

\subsubsection*{আয়াত ২৫ — সূরা আল-বাকারা (২) — আয়াত ১৫২}

\subsubsection*{১. সূত্র কার্ড}

\begin{table}[h]\centering\small
\begin{tabular}{ll}
বিষয় & বিবরণ \\
\hline
\textbf{সূরা} & আল-বাকারা (২) — মাদানী \\
\textbf{আয়াত} & ১৫২ \\
\textbf{মূল বিষয়} & স্মরণের প্রতিদান — স্মরণ; আল্লাহ বান্দার সাথে থাকার ঘোষণা \\
\textbf{প্রসঙ্গ} & হজ ও ইবাদতের পর আল্লাহর স্মরণের নির্দেশ \\
\hline
\end{tabular}
\end{table}

\subsubsection*{২. মূল আরবি}

\begin{center}\arab{فَاذْكُرُونِي أَذْكُرْكُمْ وَاشْكُرُوا لِي وَلَا تَكْفُرُونِ}\end{center}
\textbf{উচ্চারণ:} ফাযকুরূনী আজকুরকুম ওয়াশকুরূ লী ওয়ালা- তাকফুরূন।

\subsubsection*{৩. বাংলা অর্থ (\lat{pure})}

\textbf{\lat{“}অতএব তোমরা আমাকে স্মরণ করো, আমি তোমাদের স্মরণ করব। আর আমার শোকর করো এবং আমার প্রতি অকৃতজ্ঞ হয়ো না।\lat{”}}

\subsubsection*{৪. সংক্ষিপ্ত তাফসীর}

\begin{itemize}
\item \textbf{ইবনে কাসীর:} স্মরণের প্রতিদান — আল্লাহর স্মরণ; পাপী ব্যক্তির জন্য এ আয়াতের অর্থ — তাওবা ও ইস্তিগফারের মাধ্যমে আল্লাহকে স্মরণ করলেই আল্লাহ তাকে স্মরণ করেন, অর্থাৎ রহমতের দৃষ্টিতে দেখেন।
\item \textbf{সা'দী:} \lat{“}আমি তোমাদের স্মরণ করব\lat{”} — ক্ষমা, রহমত ও সাহায্যের সাথে; এটাই পাপীর ফিরে আসার সবচেয়ে সুন্দর ফল।
\end{itemize}
\subsubsection*{৫. সংশ্লিষ্ট দলিল}

\begin{itemize}
\item \textbf{হাদিস (মুসলিম ২৬৭৫):} \lat{“}আমি আমার বান্দার ধারণা অনুযায়ী — আর সে আমাকে স্মরণ করলে আমি তার সাথে থাকি।\lat{”}
\end{itemize}
\medskip\hrule\medskip

\subsection*{থিম ৭ — নবীর ফিরে আসা: ইউনুস (আঃ) (\lat{The} \lat{Prophet's} \lat{Return})}

\subsubsection*{আয়াত ২৬ — সূরা আল-আম্বিয়া (২১) — আয়াত ৮৭-৮৮}

\subsubsection*{১. সূত্র কার্ড}

\begin{table}[h]\centering\small
\begin{tabular}{ll}
বিষয় & বিবরণ \\
\hline
\textbf{সূরা} & আল-আম্বিয়া (২১) — মক্কী \\
\textbf{আয়াত} & ৮৭-৮৮ \\
\textbf{মূল বিষয়} & নবীও ভুল করেন, স্বীকার করেন, আল্লাহ তাকে উদ্ধার করেন — পাপীর দোয়ার মডেল \\
\textbf{প্রসঙ্গ} & ইউনুস (আঃ)-এর মাছের পেটের ঘটনা \\
\hline
\end{tabular}
\end{table}

\subsubsection*{২. মূল আরবি}

\begin{center}\arab{وَذَا النُّونِ إِذْ ذَهَبَ مُغَاضِبًا فَظَنَّ أَنْ لَنْ نَقْدِرَ عَلَيْهِ فَنَادَىٰ فِي الظُّلُمَاتِ أَنْ لَا إِلَٰهَ إِلَّا أَنْتَ سُبْحَانَكَ إِنِّي كُنْتُ مِنَ الظَّالِمِينَ}\end{center}
\textbf{উচ্চারণ:} ওয়া যান নূনি ইয যাহাবা মুগা-দিবান ফাজান্না আল লান নাকদিরা 'আলাইহি ফানা-দা- ফিজ যুলুমা-তি আল লা- ইলা-হা ইল্লা- আনতা সুবহা-নাকা ইন্নী কুনতু মিনাজ যা-লিমীন।

\begin{center}\arab{فَاسْتَجَبْنَا لَهُ وَنَجَّيْنَاهُ مِنَ الْغَمِّ ۚ وَكَذَٰلِكَ نُنْجِي الْمُؤْمِنِينَ}\end{center}
\textbf{উচ্চারণ:} ফাসতাজাবনা- লাহূ ওয়া নাজ্জাইনা-হু মিনাল গাম্ম। ওয়া কাযা-লিকা নুনজিল মু'মিনীন।

\subsubsection*{৩. বাংলা অর্থ (\lat{pure})}

\textbf{\lat{“}আর ইউনুসের কথা স্মরণ করো, যখন সে রাগান্বিত হয়ে চলে গিয়েছিল এবং মনে করেছিল যে আমরা তার উপর কখনো সংকীর্ণ করব না (ক্ষমতার বাইরে নয়)। অতঃপর সে অন্ধকারের মধ্যে ডেকে বলল: তুমি ছাড়া কোনো সত্য ইলাহ নেই — তুমি পবিত্র; নিশ্চয়ই আমি জালিমদের অন্তর্ভুক্ত।\lat{”}} (২১:৮৭)

\textbf{\lat{“}অতঃপর আমি তার দোয়ার জবাব দিলাম এবং তাকে দুশ্চিন্তা থেকে উদ্ধার করলাম। আর এভাবেই আমি মুমিনদেরকে উদ্ধার করি।\lat{”}} (২১:৮৮)

\subsubsection*{৪. শব্দের ব্যাখ্যা}

\begin{table}[h]\centering\small
\begin{tabular}{ll}
শব্দ & অর্থ \\
\hline
\textbf{যান নূন (\arab{ذو} \arab{النون})} & মাছওয়ালা — ইউনুস (আঃ) \\
\textbf{যুলুমাত (\arab{ظلمات})} & অন্ধকার — মাছের পেট, সমুদ্র ও রাতের অন্ধকার \\
\textbf{আন্নী কুনতু মিনাজ জালিমীন} & আমি জালিমদের অন্তর্ভুক্ত — নিজেকে দোষী বলা \\
\textbf{নুনজিল মু'মিনীন} & মুমিনদের উদ্ধার করি — প্রতিশ্রুতি \\
\hline
\end{tabular}
\end{table}

\subsubsection*{৫. সংক্ষিপ্ত তাফসীর}

\begin{itemize}
\item \textbf{ইবনে কাসীর:} ইউনুস (আঃ) জাতির উপর রাগ করে চলে গিয়েছিলেন — তারপর তিন অন্ধকারে ডেকে এই দোয়া করলেন; আল্লাহ তা কবুল করলেন এবং \lat{“}এভাবেই আমি মুমিনদের উদ্ধার করি\lat{”} বললেন। শিক্ষা — পাপী ব্যক্তির জন্য মডেল দোয়া এটি: প্রথমে আল্লাহর প্রশংসা, তারপর নিজের দোষ স্বীকার — কোনো অজুহাত নেই।
\item \textbf{সা'দী:} \lat{“}আমি জালিমদের অন্তর্ভুক্ত\lat{”} — নিজেকে দোষী মানার মধ্যেই উদ্ধারের চাবি; আল্লাহর রহমত এমন — দোষ স্বীকারকারীকে তিনি উদ্ধার করেন।
\item \textbf{মাআরিফুল কুরআন:} নবীজি (সা.) বলেছেন — কোনো মুসলিম এই দোয়া দিয়ে দোয়া করলে আল্লাহ সাড়া দেন (তিরমিযী ৩৫০৫)। অর্থ — এই দোয়া প্রতিটি পাপীর ভাষা; কুরআন ও হাদিস দুই-ই এটি শিখিয়েছে।
\end{itemize}
\subsubsection*{৬. সংশ্লিষ্ট দলিল}

\begin{itemize}
\item \textbf{হাদিস (তিরমিযী ৩৫০৫):} ইউনুসের দোয়ার ফজিলত — হাসান-সহীহ।
\item \textbf{আয়াত (সাফফাত ৩৭:১৪২-১৪৪):} ঘটনার বিস্তারিত।
\end{itemize}
\medskip\hrule\medskip

\subsection*{থিম ৮ — তাওবার দৃষ্টান্ত ও সীমারেখা (\lat{Examples} \lat{\&} \lat{Boundaries})}

\subsubsection*{আয়াত ২৭ — সূরা আত-তাওবা (৯) — আয়াত ১১৮}

\subsubsection*{১. সূত্র কার্ড}

\begin{table}[h]\centering\small
\begin{tabular}{ll}
বিষয় & বিবরণ \\
\hline
\textbf{সূরা} & আত-তাওবা (৯) — মাদানী \\
\textbf{আয়াত} & ১১৮ \\
\textbf{মূল বিষয়} & তিন সাহাবির তাওবা কবুল — দুনিয়ার সবচেয়ে বড় তাওবার উদাহরণ \\
\textbf{প্রসঙ্গ} & তাবুকের যুদ্ধে পিছিয়ে পড়া কাব, মুরারা ও হিলাল (রাঃ)-এর ঘটনা \\
\hline
\end{tabular}
\end{table}

\subsubsection*{২. মূল আরবি}

\begin{center}\arab{وَعَلَى الثَّلَاثَةِ الَّذِينَ خُلِّفُوا حَتَّىٰ إِذَا ضَاقَتْ عَلَيْهِمُ الْأَرْضُ بِمَا رَحُبَتْ وَضَاقَتْ عَلَيْهِمْ أَنْفُسُهُمْ وَظَنُّوا أَنْ لَا مَلْجَأَ مِنَ اللَّهِ إِلَّا إِلَيْهِ ثُمَّ تَابَ عَلَيْهِمْ لِيَتُوبُوا ۚ إِنَّ اللَّهَ هُوَ التَّوَّابُ الرَّحِيمُ}\end{center}
\textbf{উচ্চারণ:} ওয়া 'আলাস ছালা-ছাতিল্লাযীনা খুল্লিফূ, হাত্তা- ইযা- দা-কাত 'আলাইহিমুল আরদু বিমা- রাহুবাত ওয়া দা-কাত 'আলাইহিম আনফুসুহুম, ওয়া জান্নু আল লা- মালজাআ মিনাল্লাহি ইল্লা- ইলাইহি, ছুম্মা তা-বা 'আলাইহিম লিয়াতূবূ। ইন্নাল্লা-হা হুওয়াত তাওওয়াবুর রাহীম।

\subsubsection*{৩. বাংলা অর্থ (\lat{pure})}

\textbf{\lat{“}আর (তিন ব্যক্তির প্রতি) যাদের ব্যাপারে সিদ্ধান্ত স্থগিত রাখা হয়েছিল — এমনকি পৃথিবী প্রশস্ত হওয়া সত্ত্বেও তাদের কাছে সংকীর্ণ হয়ে গেল, তাদের জীবন তাদের কাছে দুর্বিষহ হয়ে গেল এবং তারা নিশ্চিত হলো — আল্লাহ ছাড়া আর কোনো আশ্রয় নেই। অতঃপর তিনি তাদের প্রতি ফিরলেন, যাতে তারা তাওবা করে। নিশ্চয়ই আল্লাহ তাওবা গ্রহণকারী, পরম দয়ালু।\lat{”}}

\subsubsection*{৪. শব্দের ব্যাখ্যা}

\begin{table}[h]\centering\small
\begin{tabular}{ll}
শব্দ & অর্থ \\
\hline
\textbf{খুল্লিফূ (\arab{خلّفوا})} & পিছনে রাখা হয়েছিল — সিদ্ধান্ত স্থগিত \\
\textbf{দাকাত 'আলাইহিমুল আরদ} & পৃথিবী তাদের কাছে সংকীর্ণ হয়ে গেল \\
\textbf{জান্নু আল লা মালজা} & নিশ্চিত হলো — আল্লাহ ছাড়া আশ্রয় নেই \\
\textbf{থুম্মা তাবা 'আলাইহিম} & অতঃপর তিনি তাদের প্রতি ফিরলেন \\
\hline
\end{tabular}
\end{table}

\subsubsection*{৫. সংক্ষিপ্ত তাফসীর}

\begin{itemize}
\item \textbf{ইবনে কাসীর:} এই আয়াত কাব ইবনে মালিক, মুরারা ইবনে রাবি' ও হিলাল ইবনে উমাইয়া (রাঃ)-এর ব্যাপারে — যারা তাবুক থেকে পিছিয়ে পড়েছিল; ৫০ দিন সমাজ থেকে বিচ্ছিন্ন ছিলেন; অন্তরের সেই চাপের পর আল্লাহ তাদের তাওবা কবুল করেন। ইবনে কাসীর বলেন — এর চেয়ে বড় সুসংবাদ আর নেই: আল্লাহ স্বয়ং ঘোষণা দিলেন — \lat{“}তিনি তাদের প্রতি ফিরলেন\lat{”}।
\item \textbf{সা'দী:} \lat{“}যাতে তারা তাওবা করে\lat{”} — আল্লাহর তাওবার দিকে ফেরানোই তাওবার পূর্বশর্ত; অর্থাৎ বান্দার ফেরার ইচ্ছা নিজেও আল্লাহর দান। এখান থেকে শিক্ষা — পাপের বোঝা যত ভারীই হোক, অন্তরের আন্তরিক ফেরাই আল্লাহর রহমত টানে।
\end{itemize}
\subsubsection*{৬. সংশ্লিষ্ট দলিল}

\begin{itemize}
\item \textbf{ঘটনা:} কাব ইবনে মালিক (রাঃ) — বুখারি ৪৪১৮ (সাহাবিদের ঘটনা ফাইল ১)।
\item \textbf{আয়াত (তাওবা ৯:১০২):} পাপ স্বীকারকারীদের আশা।
\end{itemize}
\medskip\hrule\medskip

\subsubsection*{আয়াত ২৮ — সূরা আত-তাহরীম (৬৬) — আয়াত ৮}

\subsubsection*{১. সূত্র কার্ড}

\begin{table}[h]\centering\small
\begin{tabular}{ll}
বিষয় & বিবরণ \\
\hline
\textbf{সূরা} & আত-তাহরীম (৬৬) — মাদানী \\
\textbf{আয়াত} & ৮ \\
\textbf{মূল বিষয়} & তাওবায়ে নসূহ — খাঁটি তাওবার আদেশ \\
\textbf{প্রসঙ্গ} & ঈমানদারদের সরাসরি নির্দেশ \\
\hline
\end{tabular}
\end{table}

\subsubsection*{২. মূল আরবি}

\begin{center}\arab{يَا أَيُّهَا الَّذِينَ آمَنُوا تُوبُوا إِلَى اللَّهِ تَوْبَةً نَصُوحًا ۖ عَسَىٰ رَبُّكُمْ أَنْ يُكَفِّرَ عَنْكُمْ سَيِّئَاتِكُمْ وَيُدْخِلَكُمْ جَنَّاتٍ تَجْرِي مِنْ تَحْتِهَا الْأَنْهَارُ}\end{center}
\textbf{উচ্চারণ:} ইয়া- আয়্যুহাল্লাযীনা আ-মানূ তূবূ ইলাল্লা-হি তাওবাতান নাসূহা। 'আসা- রব্বুকুম আই ইউকাফফিরা 'আনকুম সাইয়ি-আতিকুম ওয়া ইউদখিলাকুম জান্না-তিন তাজরী মিন তাহতিহাল আনহা-রু।

\subsubsection*{৩. বাংলা অর্থ (\lat{pure})}

\textbf{\lat{“}হে ঈমানদারগণ! তোমরা আল্লাহর কাছে তাওবা করো — খাঁটি তাওবা (তাওবাতুন নাসূহ)। আশা করা যায় তোমাদের রব তোমাদের পাপ মোচন করবেন এবং তোমাদেরকে এমন জান্নাতে প্রবেশ করাবেন যার নিচে নদী প্রবাহিত।\lat{”}}

\subsubsection*{৪. শব্দের ব্যাখ্যা}

\begin{table}[h]\centering\small
\begin{tabular}{ll}
শব্দ & অর্থ \\
\hline
\textbf{তাওবাতান নাসূহা (\arab{توبة} \arab{نصوحا})} & খাঁটি/বিশুদ্ধ তাওবা — এমন তাওবা যা পাপ ত্যাগে পূর্ণ ও ফিরে না যাওয়ার সংকল্পে দৃঢ় \\
\hline
\end{tabular}
\end{table}

\subsubsection*{৫. সংক্ষিপ্ত তাফসীর}

\begin{itemize}
\item \textbf{ইবনে কাসীর:} উমার (রাঃ) বলেন — নসূহ তাওবা: পাপে ফিরে না যাওয়ার সংকল্প, পাপের অনুতাপ ও আল্লাহর কাছে ক্ষমা প্রার্থনা। আলী (রাঃ) বলেন — ছয়টি শর্ত: অতীতের জন্য আফসোস, কাযা/হক আদায়, কারো হক থাকলে ফিরিয়ে দেওয়া, হারাম থেকে তাওবা, আল্লাহর আনুগত্যের দৃঢ় সংকল্প ও দুধ খাওয়ানো শিশুর মতো পাপ থেকে বিচ্ছিন্ন হওয়া। (তাফসীর থেকে সংক্ষিপ্ত সারমর্ম।)
\item \textbf{সা'দী:} \lat{“}আশা করা যায়\lat{”} — এখানে আশার সাথে ক্ষমার প্রতিশ্রুতি; খাঁটি তাওবার ফল — পাপ মোচন ও জান্নাত।
\end{itemize}
\subsubsection*{৬. সংশ্লিষ্ট দলিল}

\begin{itemize}
\item \textbf{আয়াত (ফুরকান ২৫:৭০):} পাপের স্থানে নেকি।
\item \textbf{ব্যবহারিক গাইড:} তাওবার শর্ত ও ধাপ — গাইড ফাইল।
\end{itemize}
\medskip\hrule\medskip

\subsubsection*{আয়াত ২৯ — সূরা আন-নিসা (৪) — আয়াত ৪৮ (সীমারেখা — সতর্কতা)}

\subsubsection*{১. সূত্র কার্ড}

\begin{table}[h]\centering\small
\begin{tabular}{ll}
বিষয় & বিবরণ \\
\hline
\textbf{সূরা} & আন-নিসা (৪) — মাদানী \\
\textbf{আয়াত} & ৪৮ \\
\textbf{মূল বিষয়} & শিরকের উপর মৃত্যু ক্ষমার বাইরে — সীমারেখার ঘোষণা \\
\textbf{প্রসঙ্গ} & ক্ষমার প্রশস্ততার সাথে সতর্কতার ভারসাম্য \\
\hline
\end{tabular}
\end{table}

\subsubsection*{২. মূল আরবি}

\begin{center}\arab{إِنَّ اللَّهَ لَا يَغْفِرُ أَنْ يُشْرَكَ بِهِ وَيَغْفِرُ مَا دُونَ ذَٰلِكَ لِمَنْ يَشَاءُ ۚ وَمَنْ يُشْرِكْ بِاللَّهِ فَقَدِ افْتَرَىٰ إِثْمًا عَظِيمًا}\end{center}
\textbf{উচ্চারণ:} ইন্নাল্লা-হা লা- ইয়াগফিরু আই ইউশরাকা বিহী ওয়া ইয়াগফিরু মা- দূনা যা-লিকা লিমাই ইয়াশা-উ। ওয়া মাই ইউশরিক বিল্লা-হি ফাকাদিফতার ইসমান 'আযীমা।

\subsubsection*{৩. বাংলা অর্থ (\lat{pure})}

\textbf{\lat{“}নিশ্চয়ই আল্লাহ শিরক ক্ষমা করবেন না; আর এ ছাড়া অন্য যা কিছু আছে তা তিনি যাকে ইচ্ছা ক্ষমা করে দেন। আর যে আল্লাহর সাথে শিরক করে, সে তো বিশাল অপবাদ রচনা করে।\lat{”}}

\subsubsection*{৪. সংক্ষিপ্ত তাফসীর}

\begin{itemize}
\item \textbf{ইবনে কাসীর:} শিরকের উপর মৃত্যু হলে ক্ষমা নেই; শিরক ছাড়া অন্য পাপ — বান্দা তাওবা করলে ক্ষমা নিশ্চিত, তাওবা না করলে আল্লাহর ইচ্ছার উপর। অর্থ — আশার সাথে সীমারেখাও স্পষ্ট: তাওবাহীন অবস্থায় শিরকে মৃত্যু হলে দরজা বন্ধ।
\item \textbf{সা'দী:} এই আয়াত দুই দিক দেখায় — (১) ক্ষমার ব্যাপকতা: শিরক ছাড়া সব পাপই ক্ষমার আওতায়; (২) সতর্কতা: শিরকই একমাত্র অক্ষম্য অবস্থা (মৃত্যুর সময় শিরকে থাকলে)।
\end{itemize}
\subsubsection*{৫. সংশ্লিষ্ট দলিল}

\begin{itemize}
\item \textbf{আয়াত (যুমার ৩৯:৫৩):} ক্ষমার ঘোষণা — তবে তাওবার শর্তে (শিরক ত্যাগে)।
\item \textbf{হাদিস (তিরমিযী ৩৫৪০):} \lat{“}...যদি তুমি আমার সাথে শিরক না করে থাকো।\lat{”}
\end{itemize}
\medskip\hrule\medskip

\subsubsection*{আয়াত ৩০ — সূরা আন-নিসা (৪) — আয়াত ৯৩ (সতর্কতা — ভারসাম্য)}

\subsubsection*{১. সূত্র কার্ড}

\begin{table}[h]\centering\small
\begin{tabular}{ll}
বিষয় & বিবরণ \\
\hline
\textbf{সূরা} & আন-নিসা (৪) — মাদানী \\
\textbf{আয়াত} & ৯৩ \\
\textbf{মূল বিষয়} & ইচ্ছাকৃত মুমিন হত্যার কঠোর সতর্কবাণী — তাওবার বিষয়ে আলেমদের মতভেদ \\
\textbf{প্রসঙ্গ} & সবচেয়ে ভারী পাপের সতর্কবাণী \\
\hline
\end{tabular}
\end{table}

\subsubsection*{২. মূল আরবি}

\begin{center}\arab{وَمَنْ يَقْتُلْ مُؤْمِنًا مُتَعَمِّدًا فَجَزَاؤُهُ جَهَنَّمُ خَالِدًا فِيهَا وَغَضِبَ اللَّهُ عَلَيْهِ وَلَعَنَهُ وَأَعَدَّ لَهُ عَذَابًا عَظِيمًا}\end{center}
\textbf{উচ্চারণ:} ওয়া মাই ইয়াকতুল মু'মিনান মুতা'আম্মিদান ফাজাআউহু জাহান্নামু খা-লিদান ফীহা- ওয়া গাদিবাল্লা-হু 'আলাইহি ওয়া লা'আনাহু ওয়া আ'আদ্দা লাহু 'আযা-বান 'আযীমা।

\subsubsection*{৩. বাংলা অর্থ (\lat{pure})}

\textbf{\lat{“}আর যে ব্যক্তি ইচ্ছাকৃতভাবে কোনো মুমিনকে হত্যা করবে — তার প্রতিদান জাহান্নাম; সে সেখানে স্থায়ী হবে। আল্লাহ তার উপর রাগান্বিত হবেন, তাকে অভিশাপ দেবেন এবং তার জন্য প্রস্তুত রেখেছেন মহাশাস্তি।\lat{”}}

\subsubsection*{৪. সংক্ষিপ্ত তাফসীর (তাওবার প্রশ্নে আলেমদের অবস্থান)}

\begin{itemize}
\item \textbf{সতর্কবাণী:} এটি কুরআনের সবচেয়ে কঠিন সতর্কবাণীগুলোর একটি — ইচ্ছাকৃত মুমিন হত্যার ভয়াবহতা বোঝাতে।
\item \textbf{ইবনে আব্বাসের (রাঃ) মত:} তিনি এই আয়াতের ব্যাখ্যায় বলেছেন — \lat{“}তার তাওবা কবুল হবে না\lat{”} (সহীহ বুখারি ৪৭৬২)। অর্থ — এত বড় পাপ যে সাধারণ তাওবার তুলনায় এর অবস্থা ভয়াবহ।
\item \textbf{জমহুরের (সংখ্যাগরিষ্ঠ) মত:} ইবনে মাসউদ, হাসান বসরি প্রমুখ ও বহু মুফাসসির বলেন — তাওবা সব পাপের জন্যই কবুল (শিরকে মৃত্যু ছাড়া), দলিল — ফুরকান ২৫:৭০, যুমার ৩৯:৫৩ ও ১০০ জন হত্যাকারীর হাদিস (মুসলিম ২৭৬৬)। ইবনে কাসীর ও পরবর্তী অনেক মুফাসসির এই আয়াতের সতর্কবাণীকে শাস্তির ভয়াবহতা বোঝাতে বলেছেন — যাতে হত্যাকে হালকা না নেওয়া হয়।
\item \textbf{মীমাংসা (ইবনে কাসীর, ইবনে তাইমিয়্যাহসহ):} তাওবা বান্দা ও আল্লাহর মধ্যে কবুল হতে পারে; কিন্তু নিহত ব্যক্তির হক (কিসাস) থেকে যায় — নিহতের পরিবার ক্ষমা না করলে কিয়ামতে হকের হিসাব থাকবে। অর্থাৎ — আল্লাহর হক মাফ হয় তাওবায়, মানুষের হক মাফ হয় ক্ষমা/কিসাসে।
\end{itemize}
\subsubsection*{৫. সংশ্লিষ্ট দলিল}

\begin{itemize}
\item \textbf{আয়াত (ফুরকান ২৫:৭০):} তাওবার পর পাপের স্থানে নেকি — হত্যাকারীর জন্যও আয়াতের ব্যাপকতা।
\item \textbf{হাদিস (মুসলিম ২৭৬৬):} ৯৯+১ জন হত্যাকারীর তাওবা — হত্যার পরেও তাওবার দরজা।
\end{itemize}
\medskip\hrule\medskip

\subsection*{সারসংক্ষেপ (\lat{Summary})}

কুরআনের এই আয়াতগুলো থেকে তাওবা ও রহমতের পুরো ছবি দাঁড়ায়:

\begin{enumerate}
\item \textbf{রহমত সবকিছুকে ঘিরে আছে} (৭:১৫৬, ৬:১২, ৪০:৩) — কোনো পাপ রহমতের বাইরে নয়; আল্লাহ নিজের উপর রহমত লিখে নিয়েছেন।
\item \textbf{হতাশা হারাম} (৩৯:৫৩, ১২:৮৭, ১৫:৫৬) — \lat{“}আল্লাহ আমাকে ক্ষমা করবেন না\lat{”} ভাবা কাফির/পথভ্রষ্টের বৈশিষ্ট্য।
\item \textbf{তাওবার দরজা খোলা} (৪২:২৫, ৯:১০৪, ২:১৬০, ৫:৩৯, ২০:৮২, ৪:১১০) — গ্রহণকারী আল্লাহ নিজেই; যেকোনো পাপের পর।
\item \textbf{আল্লাহ তাওবাকারীকে ভালোবাসেন} (২:২২২, ৩:১৩৫-১৩৬) — ক্ষমার ঊর্ধ্বে ভালোবাসা।
\item \textbf{তাওবা = নতুন জীবন} (২৫:৭০-৭১, ১৬:১১৯, ৭:১৫৩, ১৭:২৫, ৫৩:৩২, ১১:১১৪, ৯:১০২) — পাপের স্থানে নেকি; ভালো কাজ পাপ মুছে দেয়।
\item \textbf{আল্লাহ সাহায্য ও হেদায়েত দেন} (২৯:৬৯, ২:১৮৬, ৪০:৬০, ২:১৫২) — ফিরে আসার চেষ্টা আল্লাহর সাহায্য টানে; ডাকলে সাড়া মেলে।
\item \textbf{নবীর মডেল} (২১:৮৭-৮৮) — ইউনুস (আঃ): দোষ স্বীকার $\rightarrow$ উদ্ধার।
\item \textbf{দৃষ্টান্ত ও সীমারেখা} (৯:১১৮, ৬৬:৮, ৪:৪৮, ৪:৯৩) — তিন সাহাবির তাওবা কবুল; খাঁটি তাওবার আদেশ; শিরকে মৃত্যু ও মুমিন হত্যার সতর্কবাণী — আশা ও ভয়ের ভারসাম্য।
\end{enumerate}
\begin{quote}
\textbf{মূল বার্তা:} কুরআন পাপীকে কখনো হতাশ হতে বলে না — বলে \lat{“}ফিরে এসো\lat{”}। তাওবার দরজা খোলা, গ্রহণকারী আল্লাহ নিজে, আর ফিরে এলে পুরনো পাপের বদলে নতুন জীবন লেখা হয়। সাথে সীমারেখাও আছে — শিরকে মৃত্যু, মানুষের হক, পাপে অটল থাকা — যেন রহমতের আশায় পাপে নির্ভীক না হওয়া যায়।
\end{quote}


\end{document}
